\documentclass[../main.tex]{subfiles}
\begin{document}

\chapter{Advanced GPU Programming}
\lessoninfo{
  video={M5, lesson 1 (Advanced GPU Programming)},
  sources={\cite{harris2013sharedmem} (required); \cite{kipfer2005gpusort} (required, concept level); \cite{harris2012datatransfer}, \cite{harris2015streams}, \cite{nvidia2026nsightsystems}, \cite{nvidia2026nsightcompute} (optional)},
  textbook={A18 \cite{aamodt2018gpu} ch.~3; BPG \cite{nvidia2026cudabestpractices} §9 Performance Metrics, §10 Memory Optimizations, §11 Execution Configuration Optimizations; K12 \cite{kim2012gpgpu} ch.~3, §4.1; CPG \cite{nvidia2024cudaguide} §5 Performance Guidelines; PMPP \cite{hwu2022pmpp} ch.~6},
  keywords={profiling, hotspot, data transfer, pinned memory, asynchronous copy, CUDA stream, zero copy, UVA, effective bandwidth, intrinsic function, predicated execution, warp shuffle, warp vote, cooperative groups, independent thread scheduling, cuBLAS, sorting network},
}

Lesson~3 introduced the CUDA programming model and Lesson~4 the hardware
that runs it: the multithreaded SM, its register and shared-memory
banks, the pipeline and the coalescing of global memory accesses.  This
lesson turns that knowledge into a method for making a CUDA program
fast.  It starts from the question of which programs suit a GPU at all,
and from the profiler that shows where a program spends its time.  The
time of a GPU program has three parts, data transfer between host and
device, memory access on the device, and computation, and the lesson
gives a checklist for each: pinned memory, asynchronous copies and
streams for the transfers; coalescing, alignment and bank conflicts for
the memory accesses; cheaper instructions, lower precision, intrinsics,
fewer branches and fewer atomics for the computation.  It then surveys
the newer CUDA features that remove old costs (warp shuffles and votes,
cooperative groups, independent thread scheduling, clusters, split
barriers), the libraries that should be preferred to hand-written
kernels, the other GPU programming platforms, and finally the sorting
networks of the assigned reading, the classic example of an algorithm
reshaped for the GPU.  Lessons 6 and~7 go below this level, into the
hardware mechanisms behind divergence, register files and warp
scheduling.

\section{What suits a GPU, and where the time goes}
\label{sec:l05-profiling}

Not every program benefits from a GPU.\source{M5, GPU-friendly programs
and profiling; BPG §9.1 Timing; K12 §3.5.}  A program that will run well
on one has three properties:
\begin{itemize}
  \item \textbf{massive data parallelism}: the same operation applies to
    many independent data elements, enough to fill thousands of
    threads (Lessons 1 and~2);
  \item \textbf{little host--device communication} relative to the
    computation: the data moved over the I/O interface is a small share
    of the total cost;
  \item \textbf{coalesced global memory accesses}: neighbouring threads
    touch neighbouring addresses (Lesson~4).
\end{itemize}
A program that lacks the first property has no work for the lanes; one
that lacks the second or third is limited by a bus or by the memory
system, however fast the SMs are.\margin{Amdahl's law (Lesson~1) is the
quantitative form of the second property: the host--device transfer is
part of the serial fraction unless it can be overlapped
(\cref{sec:l05-transfer}).}

\subsection{Profiling}

Optimisation starts with measurement.

\begin{definition}[Profiling, hotspot]
  The measurement of where a program spends its time and what the
  hardware does while it runs is called \term{profiling}[プロファイリング]; a
  \term{profiler}[プロファイラ] is the tool that collects those
  measurements.  A \term{hotspot}[ホットスポット] is a part of the
  program (a loop, a function, a kernel) that accounts for a large share
  of the total time and is therefore the part worth optimising.
\end{definition}

Profiling comes first in two senses.  Before any code is ported, the
\emph{CPU} program is profiled, to find the hotspots that a GPU can
accelerate: only the data-parallel hotspots are moved, and the rest stays
on the host.  Once a kernel runs on the device, the GPU profiler reports
the metrics that locate the bottleneck inside it
(\cref{tab:l05-metrics}).\tip{Measure before and after every change.
Kernel launches and asynchronous copies return to the host before the
work is done, so a CPU timer must be bracketed by
\texttt{cudaDeviceSynchronize()}, or the time measured on the device
with CUDA events (\texttt{cudaEventElapsedTime}); BPG §9.1 Timing.}

\begin{table}[tb]
  \centering
  \small
  \caption{Profiler metrics named in this lesson, what a bad value
    indicates, and where this book treats the remedy.}
  \label{tab:l05-metrics}
  \begin{tabular}{>{\raggedright\arraybackslash}p{28mm}
                  >{\raggedright\arraybackslash}p{52mm}
                  >{\raggedright\arraybackslash}p{26mm}}
    \toprule
    Metric & Bad value means & Remedy \\
    \midrule
    throughput (instructions, FLOP/s) & the SMs are idle: too little
      parallelism, or stalls on memory & \cref{sec:l05-compute}, Lesson~3 \\
    divergent branches & lanes of a warp follow different paths and
      execute one after the other & \cref{sec:l05-compute}, Lesson~6 \\
    uncoalesced (\enquote{divergent}) memory accesses & one warp load
      becomes many memory transactions & \cref{sec:l05-memory}, Lesson~4 \\
    occupancy & too few resident warps to hide latency & Lesson~3 \\
    bandwidth utilisation & achieved bandwidth far below the peak:
      bad access pattern, or the kernel is not memory-bound &
      \cref{sec:l05-memory} \\
    \bottomrule
  \end{tabular}
\end{table}

\begin{definition}[Nsight Systems, Nsight Compute]
  NVIDIA's profilers are \term{Nsight Systems}[Nsight Systems], which
  records a timeline of the whole application (CPU threads, CUDA API
  calls, memory copies and kernels on every GPU) and so shows the
  transfer and overlap problems of \cref{sec:l05-transfer}
  \cite{nvidia2026nsightsystems}, and \term{Nsight Compute}[Nsight Compute],
  which profiles one kernel in detail (instruction throughput, memory
  transactions, occupancy, bank conflicts, a Roofline chart) and shows
  the problems of \cref{sec:l05-memory,sec:l05-compute}
  \cite{nvidia2026nsightcompute}.\margin{The two Nsight tools replaced
  \texttt{nvprof} and the Visual Profiler, which do not support compute
  capability 8.0 and later \cite{nvidia2026profilerguide}.}
\end{definition}

\subsection{Three parts of the execution time}

The execution time of a GPU program decomposes as
\begin{equation}
  T = T_{\text{transfer}} + T_{\text{memory}} + T_{\text{compute}},
  \label{eq:l05-time}
\end{equation}
the time to move data between host and device, the time the kernel
spends waiting for device memory, and the time it spends computing.  All
three must be attacked, and in that order of surprise: a program ported
naively is often dominated by the first term, which does not appear in
the kernel at all.  The sum reflects the naive program, in which the
three parts are serialised; the optimisations of \cref{sec:l05-transfer}
either shorten the first term or overlap it with the other two.  The
sections that follow take the three terms in
turn; a fourth rule precedes them all: use the vendor's libraries
(\cref{sec:l05-libraries}) wherever one covers the computation, because
they already contain the optimisations below.

\section{Data-transfer optimisations}
\label{sec:l05-transfer}

A discrete GPU has its own memory, and the host reaches it only through
the I/O interface (\cref{fig:l05-host-device}, left).\source{M5, data
transfer optimisations; BPG §10.1 Data Transfer Between Host and Device;
\cite{harris2012datatransfer}; \cite{harris2015streams}.}

\begin{figure}[tb]
  \centering
  % Host and device: the CPU with its host memory and the GPU with its
% device memory, joined by the I/O interface (PCIe), over which every
% explicit data transfer travels.  Redraws the top part of the outline
% figure m5l1_dataopt0.
% Usage: % Host and device: the CPU with its host memory and the GPU with its
% device memory, joined by the I/O interface (PCIe), over which every
% explicit data transfer travels.  Redraws the top part of the outline
% figure m5l1_dataopt0.
% Usage: % Host and device: the CPU with its host memory and the GPU with its
% device memory, joined by the I/O interface (PCIe), over which every
% explicit data transfer travels.  Redraws the top part of the outline
% figure m5l1_dataopt0.
% Usage: \input{figures/l05-host-device}
% Width: 48 mm, used as the left half of a body figure (with l05-pinned-staging).
\begin{tikzpicture}[x=1mm, y=1mm, font=\scriptsize\sffamily,
                    >={Stealth[length=1.4mm]}]
  \draw[fill=ln-box] (0,6) rectangle (14,16) node[midway] {CPU};
  \draw[fill=ln-accent!15] (0,0) rectangle (14,4.5)
    node[midway] {\iflangja{ホストメモリ}{host memory}};
  \draw (7,6) -- (7,4.5);
  \draw[fill=ln-box] (34,6) rectangle (48,16) node[midway] {GPU};
  \draw[fill=ln-accent!15] (34,0) rectangle (48,4.5)
    node[midway] {\iflangja{デバイスメモリ}{device memory}};
  \draw (41,6) -- (41,4.5);
  \draw[fill=white, rounded corners=1] (17,8) rectangle (31,14)
    node[midway, align=center] {\iflangja{I/O}{I/O}\\\iflangja{インタフェース}{interface}};
  \draw[<->] (14,11) -- (17,11);
  \draw[<->] (31,11) -- (34,11);
  \node[text=ln-muted] at (24,4) {PCIe};
\end{tikzpicture}

% Width: 48 mm, used as the left half of a body figure (with l05-pinned-staging).
\begin{tikzpicture}[x=1mm, y=1mm, font=\scriptsize\sffamily,
                    >={Stealth[length=1.4mm]}]
  \draw[fill=ln-box] (0,6) rectangle (14,16) node[midway] {CPU};
  \draw[fill=ln-accent!15] (0,0) rectangle (14,4.5)
    node[midway] {\iflangja{ホストメモリ}{host memory}};
  \draw (7,6) -- (7,4.5);
  \draw[fill=ln-box] (34,6) rectangle (48,16) node[midway] {GPU};
  \draw[fill=ln-accent!15] (34,0) rectangle (48,4.5)
    node[midway] {\iflangja{デバイスメモリ}{device memory}};
  \draw (41,6) -- (41,4.5);
  \draw[fill=white, rounded corners=1] (17,8) rectangle (31,14)
    node[midway, align=center] {\iflangja{I/O}{I/O}\\\iflangja{インタフェース}{interface}};
  \draw[<->] (14,11) -- (17,11);
  \draw[<->] (31,11) -- (34,11);
  \node[text=ln-muted] at (24,4) {PCIe};
\end{tikzpicture}

% Width: 48 mm, used as the left half of a body figure (with l05-pinned-staging).
\begin{tikzpicture}[x=1mm, y=1mm, font=\scriptsize\sffamily,
                    >={Stealth[length=1.4mm]}]
  \draw[fill=ln-box] (0,6) rectangle (14,16) node[midway] {CPU};
  \draw[fill=ln-accent!15] (0,0) rectangle (14,4.5)
    node[midway] {\iflangja{ホストメモリ}{host memory}};
  \draw (7,6) -- (7,4.5);
  \draw[fill=ln-box] (34,6) rectangle (48,16) node[midway] {GPU};
  \draw[fill=ln-accent!15] (34,0) rectangle (48,4.5)
    node[midway] {\iflangja{デバイスメモリ}{device memory}};
  \draw (41,6) -- (41,4.5);
  \draw[fill=white, rounded corners=1] (17,8) rectangle (31,14)
    node[midway, align=center] {\iflangja{I/O}{I/O}\\\iflangja{インタフェース}{interface}};
  \draw[<->] (14,11) -- (17,11);
  \draw[<->] (31,11) -- (34,11);
  \node[text=ln-muted] at (24,4) {PCIe};
\end{tikzpicture}
\hspace{10mm}% A host-to-device transfer from pageable memory goes through a pinned
% staging buffer (two copies); from pinned memory the DMA engine reads the
% user buffer directly (one copy).  Original, after the description in
% harris2012datatransfer.
% Usage: % A host-to-device transfer from pageable memory goes through a pinned
% staging buffer (two copies); from pinned memory the DMA engine reads the
% user buffer directly (one copy).  Original, after the description in
% harris2012datatransfer.
% Usage: % A host-to-device transfer from pageable memory goes through a pinned
% staging buffer (two copies); from pinned memory the DMA engine reads the
% user buffer directly (one copy).  Original, after the description in
% harris2012datatransfer.
% Usage: \input{figures/l05-pinned-staging}
% Width: 48 mm, used as the right half of a body figure (with l05-host-device).
\begin{tikzpicture}[x=1mm, y=1mm, font=\scriptsize\sffamily,
                    >={Stealth[length=1.4mm]}]
  % pageable path
  \node[anchor=west, font=\scriptsize\sffamily\bfseries] at (0,27) {\iflangja{ページング可能}{pageable}};
  \draw[fill=ln-box] (0,18) rectangle (13,24) node[midway, align=center] {\iflangja{ページング}{pageable}\\\iflangja{可能配列}{array}};
  \draw[fill=ln-accent!15] (17,18) rectangle (31,24) node[midway, align=center] {\iflangja{ピン留め}{pinned}\\\iflangja{中継バッファ}{staging}};
  \draw[fill=ln-box] (35,18) rectangle (48,24) node[midway, align=center] {\iflangja{デバイス}{device}\\\iflangja{メモリ}{memory}};
  \draw[->] (13,21) -- (17,21) node[midway, above, text=ln-muted] {CPU};
  \draw[->] (31,21) -- (35,21) node[midway, above, text=ln-muted] {DMA};
  % pinned path
  \node[anchor=west, font=\scriptsize\sffamily\bfseries] at (0,12) {\iflangja{ピン留め}{pinned}};
  \draw[fill=ln-accent!15] (0,3) rectangle (13,9) node[midway, align=center] {\iflangja{ピン留め}{pinned}\\\iflangja{配列}{array}};
  \draw[fill=ln-box] (35,3) rectangle (48,9) node[midway, align=center] {\iflangja{デバイス}{device}\\\iflangja{メモリ}{memory}};
  \draw[->] (13,6) -- (35,6) node[midway, above, text=ln-muted] {DMA};
  \node[text=ln-muted] at (24,0) {\iflangja{ホスト側のコピーなし}{no host-side copy}};
\end{tikzpicture}

% Width: 48 mm, used as the right half of a body figure (with l05-host-device).
\begin{tikzpicture}[x=1mm, y=1mm, font=\scriptsize\sffamily,
                    >={Stealth[length=1.4mm]}]
  % pageable path
  \node[anchor=west, font=\scriptsize\sffamily\bfseries] at (0,27) {\iflangja{ページング可能}{pageable}};
  \draw[fill=ln-box] (0,18) rectangle (13,24) node[midway, align=center] {\iflangja{ページング}{pageable}\\\iflangja{可能配列}{array}};
  \draw[fill=ln-accent!15] (17,18) rectangle (31,24) node[midway, align=center] {\iflangja{ピン留め}{pinned}\\\iflangja{中継バッファ}{staging}};
  \draw[fill=ln-box] (35,18) rectangle (48,24) node[midway, align=center] {\iflangja{デバイス}{device}\\\iflangja{メモリ}{memory}};
  \draw[->] (13,21) -- (17,21) node[midway, above, text=ln-muted] {CPU};
  \draw[->] (31,21) -- (35,21) node[midway, above, text=ln-muted] {DMA};
  % pinned path
  \node[anchor=west, font=\scriptsize\sffamily\bfseries] at (0,12) {\iflangja{ピン留め}{pinned}};
  \draw[fill=ln-accent!15] (0,3) rectangle (13,9) node[midway, align=center] {\iflangja{ピン留め}{pinned}\\\iflangja{配列}{array}};
  \draw[fill=ln-box] (35,3) rectangle (48,9) node[midway, align=center] {\iflangja{デバイス}{device}\\\iflangja{メモリ}{memory}};
  \draw[->] (13,6) -- (35,6) node[midway, above, text=ln-muted] {DMA};
  \node[text=ln-muted] at (24,0) {\iflangja{ホスト側のコピーなし}{no host-side copy}};
\end{tikzpicture}

% Width: 48 mm, used as the right half of a body figure (with l05-host-device).
\begin{tikzpicture}[x=1mm, y=1mm, font=\scriptsize\sffamily,
                    >={Stealth[length=1.4mm]}]
  % pageable path
  \node[anchor=west, font=\scriptsize\sffamily\bfseries] at (0,27) {\iflangja{ページング可能}{pageable}};
  \draw[fill=ln-box] (0,18) rectangle (13,24) node[midway, align=center] {\iflangja{ページング}{pageable}\\\iflangja{可能配列}{array}};
  \draw[fill=ln-accent!15] (17,18) rectangle (31,24) node[midway, align=center] {\iflangja{ピン留め}{pinned}\\\iflangja{中継バッファ}{staging}};
  \draw[fill=ln-box] (35,18) rectangle (48,24) node[midway, align=center] {\iflangja{デバイス}{device}\\\iflangja{メモリ}{memory}};
  \draw[->] (13,21) -- (17,21) node[midway, above, text=ln-muted] {CPU};
  \draw[->] (31,21) -- (35,21) node[midway, above, text=ln-muted] {DMA};
  % pinned path
  \node[anchor=west, font=\scriptsize\sffamily\bfseries] at (0,12) {\iflangja{ピン留め}{pinned}};
  \draw[fill=ln-accent!15] (0,3) rectangle (13,9) node[midway, align=center] {\iflangja{ピン留め}{pinned}\\\iflangja{配列}{array}};
  \draw[fill=ln-box] (35,3) rectangle (48,9) node[midway, align=center] {\iflangja{デバイス}{device}\\\iflangja{メモリ}{memory}};
  \draw[->] (13,6) -- (35,6) node[midway, above, text=ln-muted] {DMA};
  \node[text=ln-muted] at (24,0) {\iflangja{ホスト側のコピーなし}{no host-side copy}};
\end{tikzpicture}

  \caption{Left: host and device memories are separate, and every
    explicit transfer crosses the I/O interface.  Right: a copy from
    pageable host memory is two copies, one by the CPU into a pinned
    staging buffer and one DMA; from pinned memory it is one DMA.}
  \label{fig:l05-host-device}
\end{figure}

\begin{definition}[Data transfer, PCIe]
  A \term{data transfer}[データ転送] is the copy of a buffer between host
  memory and device memory (\texttt{cudaMemcpy} and its variants).  The
  interface it crosses is \term{PCI Express (PCIe)}[PCI Express]\index{PCIe|see{PCI Express (PCIe)}}, a
  serial point-to-point bus of 1 to 16 lanes, whose bandwidth is one to
  two orders of magnitude below that of the device
  memory.\sidenote{K12 §1.2 quotes \SI{4}{\giga\bit\per\second} per lane
  for PCIe Gen~2 and \SI{8}{\giga\bit\per\second} for Gen~3, i.e.\
  about \SI{8}{\giga\byte\per\second} and \SI{16}{\giga\byte\per\second}
  per direction for 16 lanes; the V100's HBM2 delivers
  \SI{898}{\giga\byte\per\second} (BPG §9.2.1).  NVLink~4 is quoted as
  $7\times$ PCIe Gen~5 \cite{nvidia2022hopper}; Lesson~9.}
\end{definition}

The naive program copies its input to the device, runs the kernel and
copies the result back, three steps in sequence
(\cref{fig:l05-transfer-timelines}, row~0).  Four techniques shorten or
hide the two copies.

\begin{figure}[tb]
  \centering
  % Four timelines of one host-to-device-to-host job, each one data-transfer
% optimisation of Lesson 5: (0) baseline copy-kernel-copy; (1) pinned
% memory shortens the copies; (2) cudaMemcpyAsync overlaps the next copy
% with the kernel; (3) streams pipeline chunked copies and kernels, drawn
% on the scale of (1) with separate copy-in and copy-out engines.
% Redraws the outline figures m5l1_dataopt0 .. m5l1_dataopt3 (timelines).
% Usage: % Four timelines of one host-to-device-to-host job, each one data-transfer
% optimisation of Lesson 5: (0) baseline copy-kernel-copy; (1) pinned
% memory shortens the copies; (2) cudaMemcpyAsync overlaps the next copy
% with the kernel; (3) streams pipeline chunked copies and kernels, drawn
% on the scale of (1) with separate copy-in and copy-out engines.
% Redraws the outline figures m5l1_dataopt0 .. m5l1_dataopt3 (timelines).
% Usage: % Four timelines of one host-to-device-to-host job, each one data-transfer
% optimisation of Lesson 5: (0) baseline copy-kernel-copy; (1) pinned
% memory shortens the copies; (2) cudaMemcpyAsync overlaps the next copy
% with the kernel; (3) streams pipeline chunked copies and kernels, drawn
% on the scale of (1) with separate copy-in and copy-out engines.
% Redraws the outline figures m5l1_dataopt0 .. m5l1_dataopt3 (timelines).
% Usage: \input{figures/l05-transfer-timelines}
% Width: body (120 mm); the drawing is about 118 mm wide.
\begin{tikzpicture}[x=1mm, y=1mm, font=\scriptsize\sffamily,
                    >={Stealth[length=1.4mm]},
                    cp/.style={fill=orange!45, draw=ln-muted},
                    kn/.style={fill=blue!30, draw=ln-muted}]
  \def\L{30}   % left edge of the bars
  % time axis
  \draw[->] (\L,2) -- (\L+86,2) node[below left, text=ln-muted] {\iflangja{時間}{time}};
  % (0) baseline
  \node[anchor=east, align=right] at (\L-2,-2) {(0) \iflangja{基本形}{baseline}};
  \fill[cp] (\L,-4.5) rectangle (\L+22,0.5) node[midway] {\iflangja{コピー}{copy in}};
  \fill[kn] (\L+22,-4.5) rectangle (\L+54,0.5) node[midway] {\iflangja{カーネル}{kernel}};
  \fill[cp] (\L+54,-4.5) rectangle (\L+76,0.5) node[midway] {\iflangja{コピー}{copy out}};
  % (1) pinned memory
  \node[anchor=east, align=right] at (\L-2,-11) {(1) \iflangja{ピン留めメモリ}{pinned memory}};
  \fill[cp] (\L,-13.5) rectangle (\L+10,-8.5) node[midway] {\iflangja{コピー}{copy}};
  \fill[kn] (\L+10,-13.5) rectangle (\L+42,-8.5) node[midway] {\iflangja{カーネル}{kernel}};
  \fill[cp] (\L+42,-13.5) rectangle (\L+52,-8.5) node[midway] {\iflangja{コピー}{copy}};
  \draw[<->, ln-muted] (\L+53,-11) -- (\L+76,-11) node[midway, above, text=ln-muted] {\iflangja{短縮}{saved}};
  % (2) overlap with cudaMemcpyAsync
  \node[anchor=east, align=right] at (\L-2,-22.5)
    {(2) \iflangja{オーバーラップ}{overlap}\\\texttt{cudaMemcpyAsync}};
  \fill[cp] (\L,-22) rectangle (\L+10,-17) node[midway] {A\,\iflangja{入}{in}};
  \fill[cp] (\L+10,-22) rectangle (\L+20,-17) node[midway] {B\,\iflangja{入}{in}};
  \fill[kn] (\L+10,-28) rectangle (\L+42,-23) node[midway] {\iflangja{カーネル}{kernel} A};
  \fill[kn] (\L+42,-28) rectangle (\L+74,-23) node[midway] {\iflangja{カーネル}{kernel} B};
  \fill[cp] (\L+42,-22) rectangle (\L+52,-17) node[midway] {A\,\iflangja{出}{out}};
  % (3) pipelined streams: same scale as (1); each chunk copy is 10/4 = 2.5,
  % each kernel chunk 32/4 = 8.  Copy-in, kernel and copy-out each have a
  % row of their own (one copy engine per direction); ends at L+37.
  \node[anchor=east, align=right] at (\L-2,-40)
    {(3) \iflangja{パイプライン}{pipeline}\\\iflangja{ストリーム}{streams}};
  \foreach \i in {0,1,2,3} {
    \fill[cp] (\L+\i*2.5,-36.5) rectangle (\L+\i*2.5+2.5,-31.5) node[midway] {\the\numexpr\i+1\relax};
    \fill[kn] (\L+2.5+\i*8,-42.5) rectangle (\L+10.5+\i*8,-37.5) node[midway] {K\the\numexpr\i+1\relax};
    \fill[cp] (\L+10.5+\i*8,-48.5) rectangle (\L+13+\i*8,-43.5) node[midway] {\the\numexpr\i+1\relax};
  }
  \draw[<->, ln-muted] (\L+38,-46) -- (\L+52,-46) node[midway, above, text=ln-muted] {\iflangja{(1) より短縮}{saved vs (1)}};
  \node[anchor=west, text=ln-muted, align=left] at (\L,-53)
    {\iflangja{(2) の上段はコピーエンジン，下段は SM．}{Row (2): upper row = copy engines, lower row = SMs.}\\
     \iflangja{(3) の 3 段はコピー（入），カーネル，コピー（出）．}{Row (3): copy in, kernels, copy out, one engine each.}\\
     \iflangja{1--4：チャンク番号}{1--4: chunk numbers}; K: \iflangja{カーネル}{kernel}};
\end{tikzpicture}

% Width: body (120 mm); the drawing is about 118 mm wide.
\begin{tikzpicture}[x=1mm, y=1mm, font=\scriptsize\sffamily,
                    >={Stealth[length=1.4mm]},
                    cp/.style={fill=orange!45, draw=ln-muted},
                    kn/.style={fill=blue!30, draw=ln-muted}]
  \def\L{30}   % left edge of the bars
  % time axis
  \draw[->] (\L,2) -- (\L+86,2) node[below left, text=ln-muted] {\iflangja{時間}{time}};
  % (0) baseline
  \node[anchor=east, align=right] at (\L-2,-2) {(0) \iflangja{基本形}{baseline}};
  \fill[cp] (\L,-4.5) rectangle (\L+22,0.5) node[midway] {\iflangja{コピー}{copy in}};
  \fill[kn] (\L+22,-4.5) rectangle (\L+54,0.5) node[midway] {\iflangja{カーネル}{kernel}};
  \fill[cp] (\L+54,-4.5) rectangle (\L+76,0.5) node[midway] {\iflangja{コピー}{copy out}};
  % (1) pinned memory
  \node[anchor=east, align=right] at (\L-2,-11) {(1) \iflangja{ピン留めメモリ}{pinned memory}};
  \fill[cp] (\L,-13.5) rectangle (\L+10,-8.5) node[midway] {\iflangja{コピー}{copy}};
  \fill[kn] (\L+10,-13.5) rectangle (\L+42,-8.5) node[midway] {\iflangja{カーネル}{kernel}};
  \fill[cp] (\L+42,-13.5) rectangle (\L+52,-8.5) node[midway] {\iflangja{コピー}{copy}};
  \draw[<->, ln-muted] (\L+53,-11) -- (\L+76,-11) node[midway, above, text=ln-muted] {\iflangja{短縮}{saved}};
  % (2) overlap with cudaMemcpyAsync
  \node[anchor=east, align=right] at (\L-2,-22.5)
    {(2) \iflangja{オーバーラップ}{overlap}\\\texttt{cudaMemcpyAsync}};
  \fill[cp] (\L,-22) rectangle (\L+10,-17) node[midway] {A\,\iflangja{入}{in}};
  \fill[cp] (\L+10,-22) rectangle (\L+20,-17) node[midway] {B\,\iflangja{入}{in}};
  \fill[kn] (\L+10,-28) rectangle (\L+42,-23) node[midway] {\iflangja{カーネル}{kernel} A};
  \fill[kn] (\L+42,-28) rectangle (\L+74,-23) node[midway] {\iflangja{カーネル}{kernel} B};
  \fill[cp] (\L+42,-22) rectangle (\L+52,-17) node[midway] {A\,\iflangja{出}{out}};
  % (3) pipelined streams: same scale as (1); each chunk copy is 10/4 = 2.5,
  % each kernel chunk 32/4 = 8.  Copy-in, kernel and copy-out each have a
  % row of their own (one copy engine per direction); ends at L+37.
  \node[anchor=east, align=right] at (\L-2,-40)
    {(3) \iflangja{パイプライン}{pipeline}\\\iflangja{ストリーム}{streams}};
  \foreach \i in {0,1,2,3} {
    \fill[cp] (\L+\i*2.5,-36.5) rectangle (\L+\i*2.5+2.5,-31.5) node[midway] {\the\numexpr\i+1\relax};
    \fill[kn] (\L+2.5+\i*8,-42.5) rectangle (\L+10.5+\i*8,-37.5) node[midway] {K\the\numexpr\i+1\relax};
    \fill[cp] (\L+10.5+\i*8,-48.5) rectangle (\L+13+\i*8,-43.5) node[midway] {\the\numexpr\i+1\relax};
  }
  \draw[<->, ln-muted] (\L+38,-46) -- (\L+52,-46) node[midway, above, text=ln-muted] {\iflangja{(1) より短縮}{saved vs (1)}};
  \node[anchor=west, text=ln-muted, align=left] at (\L,-53)
    {\iflangja{(2) の上段はコピーエンジン，下段は SM．}{Row (2): upper row = copy engines, lower row = SMs.}\\
     \iflangja{(3) の 3 段はコピー（入），カーネル，コピー（出）．}{Row (3): copy in, kernels, copy out, one engine each.}\\
     \iflangja{1--4：チャンク番号}{1--4: chunk numbers}; K: \iflangja{カーネル}{kernel}};
\end{tikzpicture}

% Width: body (120 mm); the drawing is about 118 mm wide.
\begin{tikzpicture}[x=1mm, y=1mm, font=\scriptsize\sffamily,
                    >={Stealth[length=1.4mm]},
                    cp/.style={fill=orange!45, draw=ln-muted},
                    kn/.style={fill=blue!30, draw=ln-muted}]
  \def\L{30}   % left edge of the bars
  % time axis
  \draw[->] (\L,2) -- (\L+86,2) node[below left, text=ln-muted] {\iflangja{時間}{time}};
  % (0) baseline
  \node[anchor=east, align=right] at (\L-2,-2) {(0) \iflangja{基本形}{baseline}};
  \fill[cp] (\L,-4.5) rectangle (\L+22,0.5) node[midway] {\iflangja{コピー}{copy in}};
  \fill[kn] (\L+22,-4.5) rectangle (\L+54,0.5) node[midway] {\iflangja{カーネル}{kernel}};
  \fill[cp] (\L+54,-4.5) rectangle (\L+76,0.5) node[midway] {\iflangja{コピー}{copy out}};
  % (1) pinned memory
  \node[anchor=east, align=right] at (\L-2,-11) {(1) \iflangja{ピン留めメモリ}{pinned memory}};
  \fill[cp] (\L,-13.5) rectangle (\L+10,-8.5) node[midway] {\iflangja{コピー}{copy}};
  \fill[kn] (\L+10,-13.5) rectangle (\L+42,-8.5) node[midway] {\iflangja{カーネル}{kernel}};
  \fill[cp] (\L+42,-13.5) rectangle (\L+52,-8.5) node[midway] {\iflangja{コピー}{copy}};
  \draw[<->, ln-muted] (\L+53,-11) -- (\L+76,-11) node[midway, above, text=ln-muted] {\iflangja{短縮}{saved}};
  % (2) overlap with cudaMemcpyAsync
  \node[anchor=east, align=right] at (\L-2,-22.5)
    {(2) \iflangja{オーバーラップ}{overlap}\\\texttt{cudaMemcpyAsync}};
  \fill[cp] (\L,-22) rectangle (\L+10,-17) node[midway] {A\,\iflangja{入}{in}};
  \fill[cp] (\L+10,-22) rectangle (\L+20,-17) node[midway] {B\,\iflangja{入}{in}};
  \fill[kn] (\L+10,-28) rectangle (\L+42,-23) node[midway] {\iflangja{カーネル}{kernel} A};
  \fill[kn] (\L+42,-28) rectangle (\L+74,-23) node[midway] {\iflangja{カーネル}{kernel} B};
  \fill[cp] (\L+42,-22) rectangle (\L+52,-17) node[midway] {A\,\iflangja{出}{out}};
  % (3) pipelined streams: same scale as (1); each chunk copy is 10/4 = 2.5,
  % each kernel chunk 32/4 = 8.  Copy-in, kernel and copy-out each have a
  % row of their own (one copy engine per direction); ends at L+37.
  \node[anchor=east, align=right] at (\L-2,-40)
    {(3) \iflangja{パイプライン}{pipeline}\\\iflangja{ストリーム}{streams}};
  \foreach \i in {0,1,2,3} {
    \fill[cp] (\L+\i*2.5,-36.5) rectangle (\L+\i*2.5+2.5,-31.5) node[midway] {\the\numexpr\i+1\relax};
    \fill[kn] (\L+2.5+\i*8,-42.5) rectangle (\L+10.5+\i*8,-37.5) node[midway] {K\the\numexpr\i+1\relax};
    \fill[cp] (\L+10.5+\i*8,-48.5) rectangle (\L+13+\i*8,-43.5) node[midway] {\the\numexpr\i+1\relax};
  }
  \draw[<->, ln-muted] (\L+38,-46) -- (\L+52,-46) node[midway, above, text=ln-muted] {\iflangja{(1) より短縮}{saved vs (1)}};
  \node[anchor=west, text=ln-muted, align=left] at (\L,-53)
    {\iflangja{(2) の上段はコピーエンジン，下段は SM．}{Row (2): upper row = copy engines, lower row = SMs.}\\
     \iflangja{(3) の 3 段はコピー（入），カーネル，コピー（出）．}{Row (3): copy in, kernels, copy out, one engine each.}\\
     \iflangja{1--4：チャンク番号}{1--4: chunk numbers}; K: \iflangja{カーネル}{kernel}};
\end{tikzpicture}

  \caption{One job under the four data-transfer optimisations.  (0)
    Copy in, kernel, copy out, in sequence.  (1) Pinned memory shortens
    both copies.  (2) An asynchronous copy of the next input runs on the
    copy engine while the SMs execute the current kernel.  (3) Streams
    split input, kernel and output into chunks and pipeline them, so
    that the copy engines and the SMs are busy at the same time.}
  \label{fig:l05-transfer-timelines}
\end{figure}

\subsection{Fast transfers: pinned memory}

\begin{definition}[Pinned memory]
  Host memory that the operating system has promised not to page out or
  move is called \term{pinned memory}[ピン留めメモリ] (page-locked
  memory\index{page-locked memory|see{pinned memory}}); it can be used
  by the GPU's DMA engine, which reads and writes it directly by physical
  address.  It is allocated with \texttt{cudaMallocHost()} or
  \texttt{cudaHostAlloc()} and freed with \texttt{cudaFreeHost()}
  \cite{harris2012datatransfer}.
\end{definition}

Ordinary (\emph{pageable}) host memory, from \texttt{malloc} or
\texttt{new}, cannot be the source of a DMA, because the operating system
may move or evict its pages at any time.  A \texttt{cudaMemcpy} from
pageable memory therefore first copies the data, on the CPU, into a
pinned staging buffer owned by the driver, and only then starts the DMA
(\cref{fig:l05-host-device}, right).  A copy from pinned memory skips the
staging step, which is why pinned transfers \enquote{attain the highest
bandwidth between the host and the device} (BPG §10.1.1 Pinned
Memory).\sidenote{Measured with \SI{16}{\mebi\byte} transfers
\cite{harris2012datatransfer}: on a laptop (NVS 4200M, PCIe Gen~2) pinned
memory raised host-to-device bandwidth from \num{2.31} to
\SI{5.77}{\giga\byte\per\second}, $2.5\times$; on a desktop (GTX 680) only
from \num{5.37} to \SI{6.19}{\giga\byte\per\second}, because the faster
CPU made the staging copy cheap.}

Pinned memory is a scarce resource: every pinned page is one the
operating system cannot reclaim, and \enquote{excessive use can reduce
overall system performance} (BPG §10.1.1).  The guide adds that how much
is too much cannot be known in advance; the rule is to pin the transfer
buffers, not every allocation.  Nor is pinning free: allocating pinned
memory is far slower than \texttt{malloc}, so the buffers are allocated
once, at start-up, and reused, and a buffer transferred only once rarely
pays for its pinning.

\subsection{Overlapping computation and transfer: asynchronous copies}

\begin{definition}[Asynchronous copy, overlap]
  An \term{asynchronous copy}[非同期コピー] is a transfer that returns
  control to the host as soon as it has been \emph{issued}, before the
  data has moved: \texttt{cudaMemcpyAsync()} is \enquote{a non-blocking
  variant of \texttt{cudaMemcpy()}} (BPG §10.1.2).  The
  \term{overlap of computation and transfer}[計算と転送のオーバーラップ] is
  the resulting ability of the copy engine and the SMs to work at the
  same time: while the kernel for one data set executes, the copy of the
  next data set is in flight (\cref{fig:l05-transfer-timelines},
  row~2).  The copy engine is a DMA unit separate from the SMs;
  data-centre GPUs have at least one per direction, so a host-to-device
  and a device-to-host copy can also overlap each other.
\end{definition}

The host, too, is free during an asynchronous copy, and can prepare the
next input or consume the previous output.  Overlap requires pinned
memory: an asynchronous copy from pageable memory still performs the
staging copy on the CPU before it returns, and is then largely
synchronous with respect to the host.\caution{\texttt{cudaMemcpyAsync}
with a pageable buffer compiles and runs, but gives no overlap; the
profiler timeline (Nsight Systems) shows the difference immediately.
Use pinned buffers for every asynchronous transfer.}

\subsection{Pipelining with streams: concurrent copy and execute}

Overlap becomes a pipeline when the single job is split into chunks.

\begin{definition}[CUDA stream]
  A \term{CUDA stream}[CUDAストリーム] is a sequence of device operations
  (copies, kernel launches, events) that execute in the order they were
  issued.  Operations in \emph{different} streams have no ordering
  between them and may run concurrently or interleaved
  \cite{harris2015streams}.  Every runtime call that takes a stream
  argument (\texttt{cudaMemcpyAsync}, the fourth parameter of the
  execution configuration) is placed in that stream; calls without one go
  to the \emph{default stream}.
\end{definition}

The legacy default stream synchronises with every other stream of the
device: an operation issued to it waits for all earlier work in all
streams and blocks all later work (streams created with
\texttt{cudaStreamNonBlocking} are exempt) \cite{harris2015streams}.  One stray
\texttt{cudaMemcpy} or \texttt{<<<grid, block>>>} without a stream
argument therefore serialises a whole pipeline.\caution{CUDA~7 added a
per-thread default stream (\texttt{nvcc --default-stream per-thread})
that behaves like an ordinary stream; without that flag, the default
stream is the legacy one \cite{harris2015streams}.}

\begin{definition}[Concurrent copy and execute]
  The pipelining of a job over $n$ streams is called
  \term{concurrent copy and execute}[コピーと実行の並行処理]: the input is divided into $n$
  chunks, and chunk $k$'s copy-in, kernel and copy-out are issued to
  stream $k$.  Within a stream the three steps stay in order, so each
  kernel waits for its own input; across streams the copy of chunk $k+1$
  overlaps the kernel of chunk $k$, and the copy-out of chunk $k-1$
  overlaps both (\cref{fig:l05-transfer-timelines}, row~3; BPG §10.1.2
  Asynchronous and Overlapping Transfers with Computation).
\end{definition}

\begin{lstlisting}[language=C++, numbers=none]
const int nStreams = 4;
size_t chunk = n / nStreams;  // per-chunk size
float *h_in, *h_out, *d_in, *d_out;
cudaMallocHost(&h_in,  n * sizeof(float)); // pinned
cudaMallocHost(&h_out, n * sizeof(float));
cudaMalloc(&d_in,  n * sizeof(float));
cudaMalloc(&d_out, n * sizeof(float));
cudaStream_t s[nStreams];
for (int i = 0; i < nStreams; ++i) cudaStreamCreate(&s[i]);
for (int i = 0; i < nStreams; ++i) {
  size_t off = i * chunk, bytes = chunk * sizeof(float);
  cudaMemcpyAsync(d_in + off, h_in + off, bytes,
                  cudaMemcpyHostToDevice, s[i]);
  scale<<<chunk / 256, 256, 0, s[i]>>>(d_in + off,
                                       d_out + off, chunk);
  cudaMemcpyAsync(h_out + off, d_out + off, bytes,
                  cudaMemcpyDeviceToHost, s[i]);
}
cudaDeviceSynchronize();      // wait for all
\end{lstlisting}
The loop issues twelve operations and returns at once; the device
executes them subject only to the per-stream ordering.  The
\texttt{cudaDeviceSynchronize()} at the end waits for every stream
before the host reads \texttt{h\_out}.\tip{BPG §10.1.2 gives the
estimate: with a transfer time $t_T$ and a kernel time $t_E$, the
sequential version takes $t_E + t_T$ and the staged version about
$t_E + t_T/n$ when $t_E > t_T$ (the kernel is the long pole), or
$t_T + t_E/n$ when the transfer is.}

\begin{example}[A pipelined job]
  \label{ex:l05-pipeline}
  A job copies in \SI{64}{\mebi\byte}, runs a kernel, and copies out
  \SI{64}{\mebi\byte}.  Sequentially the copy-in takes
  \SI{8}{\milli\second}, the kernel \SI{6}{\milli\second} and the
  copy-out \SI{8}{\milli\second}: \SI{22}{\milli\second} in all.  Split
  into four chunks on four streams, each chunk needs \SI{2}{\milli\second}
  in, \SI{1.5}{\milli\second} of kernel and \SI{2}{\milli\second} out; the
  device has one copy engine per direction and runs one kernel at a time.
  \Cref{tab:l05-pipeline} traces the schedule: a kernel starts when its
  input has arrived and the previous kernel has finished; a copy-out
  starts when its kernel has finished and the previous copy-out is done.
  The job completes at \SI{11.5}{\milli\second}, a speed-up of $1.9\times$,
  and the copy-in engine is the bottleneck: it is busy from $0$ to
  \SI{8}{\milli\second} without a gap, and everything after that is the
  drain of the pipeline (the last kernel and the last copy-out,
  \SI{3.5}{\milli\second}).  More chunks shorten the drain but add
  per-operation overhead.
\end{example}

\begin{table}[tb]
  \centering
  \small
  \caption{Schedule of \cref{ex:l05-pipeline} (times in
    \si{\milli\second}).  Each column is one engine; within a stream the
    three steps stay in order.}
  \label{tab:l05-pipeline}
  \begin{tabular}{lccc}
    \toprule
    Chunk (stream) & Copy in & Kernel & Copy out \\
    \midrule
    1 & 0--2 & 2--3.5 & 3.5--5.5 \\
    2 & 2--4 & 4--5.5 & 5.5--7.5 \\
    3 & 4--6 & 6--7.5 & 7.5--9.5 \\
    4 & 6--8 & 8--9.5 & 9.5--11.5 \\
    \bottomrule
  \end{tabular}
\end{table}

\subsection{Direct host memory access: zero copy and UVA}

The fourth technique removes the explicit copy altogether, by letting
the kernel read host memory.

\begin{definition}[Zero copy]
  The direct access of host memory by device threads is called
  \term{zero copy}[ゼロコピー].  The host buffer is allocated pinned and
  \emph{mapped} into the device's address space
  (\texttt{cudaHostAlloc} with \texttt{cudaHostAllocMapped}), the kernel
  receives the corresponding device pointer
  (\texttt{cudaHostGetDevicePointer}), and each load or store in the
  kernel travels over the I/O interface instead of to device
  memory (BPG §10.1.3 Zero Copy).
\end{definition}

\begin{lstlisting}[language=C++, numbers=none]
float *h_a, *d_a;
cudaHostAlloc(&h_a, bytes, cudaHostAllocMapped);
cudaHostGetDevicePointer(&d_a, h_a, 0);
kernel<<<grid, block>>>(d_a);  // loads cross the I/O bus
\end{lstlisting}

Zero copy is a clear win on an \emph{integrated} GPU, whose memory is
physically the CPU's memory: the \enquote{copy} it replaces was a copy
from one region of the same DRAM to another.  On a discrete GPU every
access pays the I/O interface's bandwidth and latency, so it helps only
when the data is read once, in a coalesced pattern, and the latency can
be hidden by other warps; data that is reused belongs in device
memory.

\begin{definition}[Unified virtual addressing]
  Under
  \term{unified virtual addressing (UVA)}[統合仮想アドレッシング（UVA）]\index{UVA|see{unified virtual addressing (UVA)}},
  host memory and the memories of all devices share one virtual
  address space: a pointer value identifies its memory by itself, the
  runtime can tell from the address which memory a pointer refers to
  (\texttt{cudaPointerGetAttributes}), and
  \texttt{cudaMemcpy} can be called with \texttt{cudaMemcpyDefault}
  without naming the direction (BPG §10.1.4 Unified Virtual
  Addressing).  The driver and runtime hide the physically separate
  memories behind one addressing scheme.\margin{UVA needs compute
  capability 2.0 and a 64-bit host; under it a pinned buffer has the
  same host and device pointer, and peer-to-peer copies between GPUs
  (Lesson~9) become possible.  BPG §10.1.4.}
\end{definition}

Ideally a unified address space would make data movement free.  It does
not: the memories are still separate, the data must still cross the
interface, and the copy happens either explicitly or on the runtime's
initiative, with its overhead.  UVA changes the \emph{naming} of memory,
not its \emph{location}.\caution{UVA is not unified memory.  UVA (2011,
CUDA~4) is one address space over separate memories; unified memory
(\texttt{cudaMallocManaged}, CUDA~6) adds automatic migration of pages
between them on demand, with page-fault costs of its own.  Lesson~7
treats unified memory.}

Not every one of the four techniques fits every system: the gain of
pinned memory depends on the host CPU, the gain of overlap on the ratio
of transfer to kernel time, the gain of zero copy on whether the GPU is
integrated.  The rule is to try each technique and measure.

\section{Memory-access optimisations}
\label{sec:l05-memory}

Once the data is on the device, the second term of \cref{eq:l05-time} is
the time the kernel waits for device memory.\source{M5, memory-access
optimisations; BPG §9.2 Bandwidth, §10.2 Device Memory Spaces;
\cite{harris2013sharedmem}.}  The checklist repeats, as rules,
the mechanisms of Lesson~4:
\begin{itemize}
  \item use the on-chip memories: shared memory for data a block
    reuses, the L1 and L2 caches for data with locality, constant
    memory for values every thread reads;\margin{Shared memory latency
    is \enquote{roughly 100× lower than uncached global memory latency}
    \cite{harris2013sharedmem}; on the Tesla and Fermi generations
    K12 Table~1.4 lists 4--32 cycles for shared memory against hundreds
    for global memory.}
  \item \textbf{coalesce} and \textbf{align} global accesses so that a
    warp's load touches as few memory transactions as possible: on
    current GPUs a transaction is a \SI{32}{\byte} sector, and a warp of
    32 threads reading consecutive 4-byte words touches four of them
    (BPG §10.2.1 Coalesced Access to Global Memory);
  \item \textbf{reduce the number of transactions}: a structure of
    arrays instead of an array of structures, data types sized to what
    is needed, loads of 8 or 16 bytes per thread where alignment allows;
  \item \textbf{compare the achieved bandwidth with the peak}
    (\cref{def:l05-bandwidth}): a gap that persists after coalescing
    means the kernel is not memory-bound and the next term of
    \cref{eq:l05-time} is the one to attack;
  \item \textbf{avoid shared-memory bank conflicts}: 32 banks of
    32-bit words, one access per bank per cycle, so a warp whose lanes
    access different addresses in the same bank serialises (lanes
    reading the same word get a broadcast; Lesson~4;
    \cite{harris2013sharedmem}).
\end{itemize}

\begin{definition}[Theoretical and effective bandwidth]
  \label{def:l05-bandwidth}
  The \term{theoretical bandwidth}[理論帯域幅] of a memory is computed
  from its hardware specification: memory clock rate $\times$ interface
  width in bytes $\times$ 2 for double data rate (BPG §9.2.1
  Theoretical Bandwidth Calculation).  The
  \term{effective bandwidth}[実効帯域幅] of a kernel is the bandwidth it
  actually achieved,
  \begin{equation}
    \text{effective bandwidth} = \frac{(B_r + B_w)/10^9}{t},
    \label{eq:l05-effective-bw}
  \end{equation}
  with $B_r$ and $B_w$ the bytes the kernel reads and writes and $t$ its
  execution time in seconds, in \si{\giga\byte\per\second} (BPG §9.2.2
  Effective Bandwidth Calculation).\margin{BPG's own instance: the V100,
  HBM2 at \SI{877}{\mega\hertz} on 4096 bits, $0.877 \times 10^9 \times
  512 \times 2 / 10^9 = \SI{898}{\giga\byte\per\second}$ (BPG §9.2.1).}
\end{definition}

Both numbers must use the same divisor: bandwidths are decimal
(\num{e9}), and dividing the bytes by $1024^3$ instead understates the
effective bandwidth by \SI{7}{\percent} (BPG §9.2.1).  The Visual
Profiler reported a \enquote{requested} throughput, corresponding to the
effective bandwidth, and an \enquote{actual} throughput that includes
the bytes wasted by uncoalesced accesses; the gap between the two is the
coalescing loss (BPG §9.2.3 Throughput Reported by Visual Profiler).  In
Nsight Compute the same loss appears as sectors per request in the
memory workload analysis: the ideal is four for 32 consecutive 4-byte
loads.

\begin{example}[Bandwidth utilisation]
  \label{ex:l05-bandwidth}
  A GPU has HBM2 at a memory clock of \SI{1.2}{\giga\hertz} on a 5120-bit
  interface.  Its theoretical bandwidth is
  \[
    1.2 \times 10^9 \times \frac{5120}{8} \times 2 \,/\, 10^9
    = \SI{1536}{\giga\byte\per\second}.
  \]
  A kernel reads an array of \SI{256}{\mebi\byte}
  ($B_r = \num{268435456}$) and writes one of the same size, and runs in
  \SI{0.5}{\milli\second}.  By \cref{eq:l05-effective-bw},
  \[
    \frac{(\num{268435456} + \num{268435456})/10^9}{0.5 \times 10^{-3}}
    = \SI{1074}{\giga\byte\per\second},
  \]
  \SI{70}{\percent} of the peak.  The kernel is memory-bound and its
  access pattern is already good; the remaining \SI{30}{\percent} is the
  usual gap between a specification and a streaming kernel, not a sign
  of uncoalesced accesses.  Had the result been \SI{300}{\giga\byte\per\second}
  with the same byte count, the first suspect would be the access
  pattern, and the sectors per request in Nsight Compute the first
  metric to read.
\end{example}

\section{Reducing computation}
\label{sec:l05-compute}

The third term of \cref{eq:l05-time} is the computation
itself.\source{M5, reducing computation; BPG §12 Instruction
Optimization, §13 Control Flow; CPG §5.4 Maximize Instruction Throughput.}

\begin{definition}[Instruction optimisation]
  The replacement of the instructions a kernel executes by fewer or
  cheaper ones with the same (or an acceptable) result is called
  \term{instruction optimization}[命令最適化]: a cheaper operation for the same function,
  a narrower data type, a hardware function unit for a software routine,
  a predicate for a branch.
\end{definition}

Six techniques reduce computation.
\begin{description}
  \item[Cheaper instructions.]  Integer division and modulo are
    multi-instruction sequences on the SM (BPG §12.1.4); 32-bit integer
    multiplication is a single instruction on Volta and later but was
    emulated on Maxwell and Pascal; shifts and masks are single
    instructions.  For a power-of-two divisor
    $n = 2^k$, \texttt{i / n} is \texttt{i >{}> k} and \texttt{i \% n} is
    \texttt{i \& (n - 1)} for non-negative \texttt{i}.\margin{The instruction-level form of strength
    reduction (Lesson~11).  The compiler does this itself for a
    compile-time constant divisor; the rewrite matters for a run-time
    power of two.}
\begin{lstlisting}[language=C++, numbers=none]
int q1 = i / n;        // division sequence
int q2 = i >> k;       // one shift, if n == 1 << k
int r  = i & (n - 1);  // i % n, for i >= 0
\end{lstlisting}
  \item[Lower precision.]  Where the application tolerates it, 32-bit
    arithmetic instead of 64-bit, and 16- or 8-bit instead of 32-bit:
    each halving of the width doubles the throughput of the arithmetic
    units and halves the bytes moved.  FP64 throughput on consumer GPUs
    is a small fraction of FP32; FP16 pairs (\texttt{half2}) execute two
    operations per instruction (BPG §12.1).  Lesson~12 treats the formats
    and what the tensor cores do with them.\caution{Lower precision is a
    numerical decision before it is a performance one: a sum of many
    FP16 terms loses digits at every addition, and a reduction over a
    million elements in FP16 is wrong, not merely approximate.  Keep
    accumulators wider than the operands.}
  \item[Hardware special functions.]  The SM's special function units
    (Lesson~4) evaluate reciprocal, reciprocal square root, $\log_2$,
    $2^x$, sine and cosine as single instructions, at reduced accuracy.
    \begin{definition}[Intrinsic function]
      An \term{intrinsic function}[組み込み関数] is a function the
      compiler translates directly into a hardware instruction rather
      than into a call to a library routine, such as
      \texttt{\_\_sinf()}, \texttt{\_\_expf()}, \texttt{\_\_fdividef()};
      the standard function \texttt{rsqrtf()} likewise maps to the
      SFU's reciprocal square root.  The reciprocal square root \enquote{should always be
      invoked explicitly as \texttt{rsqrtf()}} rather than written as
      \texttt{1.0f / sqrtf(x)}, which the compiler cannot always merge
      (BPG §12.1.6).
    \end{definition}
\begin{lstlisting}[language=C++, numbers=none]
float y1 = 1.0f / sqrtf(x);  // two ops
float y2 = rsqrtf(x);        // one SFU instruction
\end{lstlisting}
    The math libraries (\cref{sec:l05-libraries}) use the intrinsics
    automatically; so does the compiler, for a whole file, under
    \texttt{nvcc -use\_fast\_math}, which also flushes denormals to zero
    and selects the approximate division and square root
    (BPG §12.1.9).\caution{\texttt{-use\_fast\_math} changes results,
    not only speed: check the output against the precise build.}
  \item[Fewer branches.]  A branch whose condition differs between the
    lanes of a warp makes the warp execute both paths (Lesson~6).  Two
    remedies: write the condition so that whole warps take the same
    path, as in a condition on \texttt{threadIdx.x / warpSize} rather
    than on \texttt{threadIdx.x} (CPG §5.4.2 Control Flow Instructions),
    and let the compiler replace short branches by predication.
    \begin{definition}[Predicated execution]
      In \term{predicated execution}[プレディケート実行] both sides of a
      conditional are issued to every lane, and each instruction carries
      a predicate register that decides, per lane, whether its result is
      written.  No branch instruction is executed and no warp diverges;
      the cost is that every lane spends the cycles of both sides.  The
      compiler chooses predication only when the number of instructions
      controlled by the condition is below a threshold
      (BPG §13.2 Branch Predication).
    \end{definition}
    Lesson~6 shows the hardware side: the SIMT stack, the active mask,
    and the PTX guard predicates of Lesson~10.
  \item[Avoid atomics.]  An atomic read-modify-write on a location that
    many threads update serialises them at the memory or in the L2
    cache.  Atomics have become much cheaper over the GPU generations,
    but a reduction into one counter by thousands of threads is still a
    serialisation; reduce within the block first (shared memory, or the
    warp shuffles of \cref{sec:l05-features}) and issue one atomic per
    block.\margin{Kepler implemented shared-memory atomics as a
    lock/update/unlock sequence; Maxwell added native shared-memory
    atomic instructions for 32-bit integers
    \cite{nvidia2026maxwelltuning}: atomics became \enquote{much cheaper
    than before} through changes of this kind.}
  \item[Tensor cores.]  A matrix multiply-accumulate on the tensor cores
    (Lesson~1) runs an order of magnitude faster than the same
    operation on the FP32 lanes.  Any computation that can be phrased as
    a matrix product should be, and the libraries do so when the shapes
    and types allow (Lesson~12).
\end{description}

\begin{keypoint}
  The three terms of \cref{eq:l05-time} have three checklists.
  Transfers: pin, overlap, pipeline, or map.  Memory: cache, coalesce,
  align, count transactions, avoid bank conflicts, compare with the
  peak.  Computation: cheaper instructions, narrower types, intrinsics,
  warp-uniform conditions, fewer atomics, tensor cores.  Profile before
  and after each step.
\end{keypoint}

\section{Newer CUDA features}
\label{sec:l05-features}

Several costs of the checklists above have been removed by hardware and
API additions since the original CUDA model.\source{M5, newer CUDA
features; CPG §4.1 SIMT Architecture, §7.19 Warp Vote Functions,
§7.22 Warp Shuffle Functions, §7.26 Asynchronous Barrier, §8 Cooperative
Groups; \cite{nvidia2022hopper}.}

\subsection{Warp-level operations: shuffle and vote}

Registers are private to a thread.  In the original model two threads of
the same warp could exchange a value only through shared memory: a store,
a \texttt{\_\_syncthreads()}, a load, with the risk of a bank conflict
on the way.  Warp-level operations move the value directly between
lanes.

\begin{definition}[Warp shuffle]
  A \term{warp shuffle}[ワープシャッフル] is an instruction by which each
  lane of a warp reads a register of another lane of the same warp,
  with no shared memory and no barrier (CPG §7.22 Warp Shuffle
  Functions).  The four variants differ in which lane is read
  (\cref{tab:l05-shuffle}); every variant takes a \emph{mask} of the
  lanes that participate, the value, and the lane selector.
\end{definition}

\begin{table}[tb]
  \centering
  \small
  \caption{The warp shuffle functions.  \texttt{mask} names the
    participating lanes (\texttt{0xffffffff} for all 32); the optional
    \texttt{width} splits the warp into independent segments.}
  \label{tab:l05-shuffle}
  \begin{tabular}{>{\raggedright\arraybackslash}p{52mm}
                  >{\raggedright\arraybackslash}p{58mm}}
    \toprule
    Function & Lane $i$ receives the value of lane \\
    \midrule
    \texttt{\_\_shfl\_sync(mask, v, src)} & \texttt{src} (one lane
      broadcasts to all) \\
    \texttt{\_\_shfl\_up\_sync(mask, v, d)} & $i - d$ (lanes below $d$
      keep their own value) \\
    \texttt{\_\_shfl\_down\_sync(mask, v, d)} & $i + d$ (the top $d$
      lanes keep their own value) \\
    \texttt{\_\_shfl\_xor\_sync(mask, v, m)} & $i \oplus m$ (a
      butterfly exchange) \\
    \bottomrule
  \end{tabular}
\end{table}

\begin{figure}[tb]
  \centering
  % Warp shuffle: eight lanes of a warp exchange the value of R1 with their
% neighbour (lane XOR 1) in one instruction, without touching shared
% memory.  Redraws the outline figure m5l1_cudafeat.
% Usage: % Warp shuffle: eight lanes of a warp exchange the value of R1 with their
% neighbour (lane XOR 1) in one instruction, without touching shared
% memory.  Redraws the outline figure m5l1_cudafeat.
% Usage: % Warp shuffle: eight lanes of a warp exchange the value of R1 with their
% neighbour (lane XOR 1) in one instruction, without touching shared
% memory.  Redraws the outline figure m5l1_cudafeat.
% Usage: \input{figures/l05-shuffle-xor}
% Width: body (120 mm); the drawing is about 100 mm wide.
\begin{tikzpicture}[x=1mm, y=1mm, font=\scriptsize\sffamily,
                    >={Stealth[length=1.4mm]}]
  \def\cw{12}
  \node[anchor=east, text=ln-muted] at (-1,19) {\iflangja{実行前}{before}};
  \node[anchor=east, text=ln-muted] at (-1,3) {\iflangja{実行後}{after}};
  \foreach \i/\v/\c in {0/a/red!35, 1/b/blue!30, 2/c/green!40, 3/d/yellow!60,
                        4/e/purple!35, 5/f/orange!45, 6/g/cyan!35, 7/h/gray!40} {
    \node[text=ln-muted] at (\i*\cw+\cw/2,25) {\iflangja{レーン}{lane} \i};
    \fill[\c, draw=ln-muted] (\i*\cw+1.5,16) rectangle (\i*\cw+\cw-1.5,22);
    \node at (\i*\cw+\cw/2,19) {R1 = $\v$};
  }
  \foreach \i/\v/\c in {0/b/blue!30, 1/a/red!35, 2/d/yellow!60, 3/c/green!40,
                        4/f/orange!45, 5/e/purple!35, 6/h/gray!40, 7/g/cyan!35} {
    \fill[\c, draw=ln-muted] (\i*\cw+1.5,0) rectangle (\i*\cw+\cw-1.5,6);
    \node at (\i*\cw+\cw/2,3) {R1 = $\v$};
  }
  % crossing arrows for each pair
  \foreach \i in {0,2,4,6} {
    \draw[->] (\i*\cw+\cw/2,15.5) -- (\i*\cw+\cw+\cw/2,6.5);
    \draw[->] (\i*\cw+\cw+\cw/2,15.5) -- (\i*\cw+\cw/2,6.5);
  }
  \node[anchor=west, text=ln-muted] at (0,-4)
    {\texttt{\_\_shfl\_xor\_sync(0xffffffff, R1, 1)}: \iflangja{レーン $i$ はレーン $i \oplus 1$ の値を受け取る}{lane $i$ receives the value of lane $i \oplus 1$}};
\end{tikzpicture}

% Width: body (120 mm); the drawing is about 100 mm wide.
\begin{tikzpicture}[x=1mm, y=1mm, font=\scriptsize\sffamily,
                    >={Stealth[length=1.4mm]}]
  \def\cw{12}
  \node[anchor=east, text=ln-muted] at (-1,19) {\iflangja{実行前}{before}};
  \node[anchor=east, text=ln-muted] at (-1,3) {\iflangja{実行後}{after}};
  \foreach \i/\v/\c in {0/a/red!35, 1/b/blue!30, 2/c/green!40, 3/d/yellow!60,
                        4/e/purple!35, 5/f/orange!45, 6/g/cyan!35, 7/h/gray!40} {
    \node[text=ln-muted] at (\i*\cw+\cw/2,25) {\iflangja{レーン}{lane} \i};
    \fill[\c, draw=ln-muted] (\i*\cw+1.5,16) rectangle (\i*\cw+\cw-1.5,22);
    \node at (\i*\cw+\cw/2,19) {R1 = $\v$};
  }
  \foreach \i/\v/\c in {0/b/blue!30, 1/a/red!35, 2/d/yellow!60, 3/c/green!40,
                        4/f/orange!45, 5/e/purple!35, 6/h/gray!40, 7/g/cyan!35} {
    \fill[\c, draw=ln-muted] (\i*\cw+1.5,0) rectangle (\i*\cw+\cw-1.5,6);
    \node at (\i*\cw+\cw/2,3) {R1 = $\v$};
  }
  % crossing arrows for each pair
  \foreach \i in {0,2,4,6} {
    \draw[->] (\i*\cw+\cw/2,15.5) -- (\i*\cw+\cw+\cw/2,6.5);
    \draw[->] (\i*\cw+\cw+\cw/2,15.5) -- (\i*\cw+\cw/2,6.5);
  }
  \node[anchor=west, text=ln-muted] at (0,-4)
    {\texttt{\_\_shfl\_xor\_sync(0xffffffff, R1, 1)}: \iflangja{レーン $i$ はレーン $i \oplus 1$ の値を受け取る}{lane $i$ receives the value of lane $i \oplus 1$}};
\end{tikzpicture}

% Width: body (120 mm); the drawing is about 100 mm wide.
\begin{tikzpicture}[x=1mm, y=1mm, font=\scriptsize\sffamily,
                    >={Stealth[length=1.4mm]}]
  \def\cw{12}
  \node[anchor=east, text=ln-muted] at (-1,19) {\iflangja{実行前}{before}};
  \node[anchor=east, text=ln-muted] at (-1,3) {\iflangja{実行後}{after}};
  \foreach \i/\v/\c in {0/a/red!35, 1/b/blue!30, 2/c/green!40, 3/d/yellow!60,
                        4/e/purple!35, 5/f/orange!45, 6/g/cyan!35, 7/h/gray!40} {
    \node[text=ln-muted] at (\i*\cw+\cw/2,25) {\iflangja{レーン}{lane} \i};
    \fill[\c, draw=ln-muted] (\i*\cw+1.5,16) rectangle (\i*\cw+\cw-1.5,22);
    \node at (\i*\cw+\cw/2,19) {R1 = $\v$};
  }
  \foreach \i/\v/\c in {0/b/blue!30, 1/a/red!35, 2/d/yellow!60, 3/c/green!40,
                        4/f/orange!45, 5/e/purple!35, 6/h/gray!40, 7/g/cyan!35} {
    \fill[\c, draw=ln-muted] (\i*\cw+1.5,0) rectangle (\i*\cw+\cw-1.5,6);
    \node at (\i*\cw+\cw/2,3) {R1 = $\v$};
  }
  % crossing arrows for each pair
  \foreach \i in {0,2,4,6} {
    \draw[->] (\i*\cw+\cw/2,15.5) -- (\i*\cw+\cw+\cw/2,6.5);
    \draw[->] (\i*\cw+\cw+\cw/2,15.5) -- (\i*\cw+\cw/2,6.5);
  }
  \node[anchor=west, text=ln-muted] at (0,-4)
    {\texttt{\_\_shfl\_xor\_sync(0xffffffff, R1, 1)}: \iflangja{レーン $i$ はレーン $i \oplus 1$ の値を受け取る}{lane $i$ receives the value of lane $i \oplus 1$}};
\end{tikzpicture}

  \caption{\texttt{\_\_shfl\_xor\_sync} with lane mask 1: the eight
    lanes shown swap R1 with their neighbour in one instruction.}
  \label{fig:l05-shuffle-xor}
\end{figure}

\Cref{fig:l05-shuffle-xor} shows the simplest case, a pairwise swap.
The important use is the reduction: a sum over the 32 lanes of a warp
takes five shuffles, one per bit of the lane index, and every lane ends
with the total (\cref{fig:l05-butterfly}, \cref{ex:l05-butterfly}).
\Needspace{6\baselineskip}
\begin{lstlisting}[language=C++, numbers=none]
__device__ float warp_sum(float v) {
  for (int m = 16; m > 0; m >>= 1)
    v += __shfl_xor_sync(0xffffffff, v, m);
  return v;   // every lane holds the warp total
}
\end{lstlisting}
A block reduction then needs shared memory only once, for the per-warp
totals, and a grid reduction one atomic per
block.\margin{The \texttt{\_sync} forms date from CUDA~9, which
deprecated the old \texttt{\_\_shfl()} family without a mask when Volta
introduced independent thread scheduling (below); the mask tells the
hardware which lanes must have arrived before the exchange.  Compute
capability 8.0 adds warp reductions as single functions,
\texttt{\_\_reduce\_add\_sync} and relatives (CPG §7.21 Warp Reduce
Functions).}

\begin{example}[Butterfly reduction]
  \label{ex:l05-butterfly}
  Eight lanes hold the values $3, 1, 4, 1, 5, 9, 2, 6$ (sum 31), and the
  loop runs with masks $4, 2, 1$.  After the first round lane $i$ holds
  $v_i + v_{i \oplus 4}$: lanes 0--3 hold $8, 10, 6, 7$ and lanes 4--7
  the same values.  After mask 2, lanes 0 and~1 hold $8 + 6 = 14$ and
  $10 + 7 = 17$, repeated across the warp; after mask 1, every lane holds
  $14 + 17 = 31$ (\cref{fig:l05-butterfly}).  Three rounds for 8 lanes,
  $\log_2 32 = 5$ for a full warp.
\end{example}

\begin{figure}[tb]
  \centering
  % Butterfly reduction with __shfl_xor_sync over eight lanes: three rounds
% with lane masks 4, 2 and 1; after the last round every lane holds the
% sum.  Original (the values are the chapter's own example).
% Usage: % Butterfly reduction with __shfl_xor_sync over eight lanes: three rounds
% with lane masks 4, 2 and 1; after the last round every lane holds the
% sum.  Original (the values are the chapter's own example).
% Usage: % Butterfly reduction with __shfl_xor_sync over eight lanes: three rounds
% with lane masks 4, 2 and 1; after the last round every lane holds the
% sum.  Original (the values are the chapter's own example).
% Usage: \input{figures/l05-butterfly}
% Width: body (120 mm); the drawing is about 104 mm wide.
\begin{tikzpicture}[x=1mm, y=1mm, font=\scriptsize\sffamily,
                    >={Stealth[length=1.2mm]}]
  \def\cw{11}
  \def\rh{11}
  \foreach \i in {0,...,7} {
    \node[text=ln-muted] at (\i*\cw+\cw/2,3*\rh+5) {\iflangja{レーン}{lane} \i};
  }
  % rows: round 0 (input), after xor 4, after xor 2, after xor 1
  \foreach \r/\lbl/\vals in {
      3/{\iflangja{入力}{input}}/{3,1,4,1,5,9,2,6},
      2/{$\oplus 4$}/{8,10,6,7,8,10,6,7},
      1/{$\oplus 2$}/{14,17,14,17,14,17,14,17},
      0/{$\oplus 1$}/{31,31,31,31,31,31,31,31}} {
    \node[anchor=east, text=ln-muted] at (-1,\r*\rh+2.5) {\lbl};
    \foreach \v [count=\i from 0] in \vals {
      \draw[fill=ln-box] (\i*\cw+1.5,\r*\rh) rectangle (\i*\cw+\cw-1.5,\r*\rh+5);
      \node at (\i*\cw+\cw/2,\r*\rh+2.5) {\v};
    }
  }
  % arrows: xor 4 (row 3 -> 2)
  \foreach \i in {0,1,2,3} {
    \draw[->, ln-accent] (\i*\cw+\cw/2,3*\rh-0.3) -- (\i*\cw+4*\cw+\cw/2,2*\rh+5.3);
    \draw[->, ln-accent] (\i*\cw+4*\cw+\cw/2,3*\rh-0.3) -- (\i*\cw+\cw/2,2*\rh+5.3);
  }
  % xor 2 (row 2 -> 1)
  \foreach \i in {0,1,4,5} {
    \draw[->, ln-accent] (\i*\cw+\cw/2,2*\rh-0.3) -- (\i*\cw+2*\cw+\cw/2,\rh+5.3);
    \draw[->, ln-accent] (\i*\cw+2*\cw+\cw/2,2*\rh-0.3) -- (\i*\cw+\cw/2,\rh+5.3);
  }
  % xor 1 (row 1 -> 0)
  \foreach \i in {0,2,4,6} {
    \draw[->, ln-accent] (\i*\cw+\cw/2,\rh-0.3) -- (\i*\cw+\cw+\cw/2,5.3);
    \draw[->, ln-accent] (\i*\cw+\cw+\cw/2,\rh-0.3) -- (\i*\cw+\cw/2,5.3);
  }
  \node[anchor=west, text=ln-muted] at (0,-4)
    {\iflangja{各段で，レーン $i$ は自分の値にレーン $i \oplus m$ の値を加える}{in each round lane $i$ adds the value of lane $i \oplus m$ to its own}};
\end{tikzpicture}

% Width: body (120 mm); the drawing is about 104 mm wide.
\begin{tikzpicture}[x=1mm, y=1mm, font=\scriptsize\sffamily,
                    >={Stealth[length=1.2mm]}]
  \def\cw{11}
  \def\rh{11}
  \foreach \i in {0,...,7} {
    \node[text=ln-muted] at (\i*\cw+\cw/2,3*\rh+5) {\iflangja{レーン}{lane} \i};
  }
  % rows: round 0 (input), after xor 4, after xor 2, after xor 1
  \foreach \r/\lbl/\vals in {
      3/{\iflangja{入力}{input}}/{3,1,4,1,5,9,2,6},
      2/{$\oplus 4$}/{8,10,6,7,8,10,6,7},
      1/{$\oplus 2$}/{14,17,14,17,14,17,14,17},
      0/{$\oplus 1$}/{31,31,31,31,31,31,31,31}} {
    \node[anchor=east, text=ln-muted] at (-1,\r*\rh+2.5) {\lbl};
    \foreach \v [count=\i from 0] in \vals {
      \draw[fill=ln-box] (\i*\cw+1.5,\r*\rh) rectangle (\i*\cw+\cw-1.5,\r*\rh+5);
      \node at (\i*\cw+\cw/2,\r*\rh+2.5) {\v};
    }
  }
  % arrows: xor 4 (row 3 -> 2)
  \foreach \i in {0,1,2,3} {
    \draw[->, ln-accent] (\i*\cw+\cw/2,3*\rh-0.3) -- (\i*\cw+4*\cw+\cw/2,2*\rh+5.3);
    \draw[->, ln-accent] (\i*\cw+4*\cw+\cw/2,3*\rh-0.3) -- (\i*\cw+\cw/2,2*\rh+5.3);
  }
  % xor 2 (row 2 -> 1)
  \foreach \i in {0,1,4,5} {
    \draw[->, ln-accent] (\i*\cw+\cw/2,2*\rh-0.3) -- (\i*\cw+2*\cw+\cw/2,\rh+5.3);
    \draw[->, ln-accent] (\i*\cw+2*\cw+\cw/2,2*\rh-0.3) -- (\i*\cw+\cw/2,\rh+5.3);
  }
  % xor 1 (row 1 -> 0)
  \foreach \i in {0,2,4,6} {
    \draw[->, ln-accent] (\i*\cw+\cw/2,\rh-0.3) -- (\i*\cw+\cw+\cw/2,5.3);
    \draw[->, ln-accent] (\i*\cw+\cw+\cw/2,\rh-0.3) -- (\i*\cw+\cw/2,5.3);
  }
  \node[anchor=west, text=ln-muted] at (0,-4)
    {\iflangja{各段で，レーン $i$ は自分の値にレーン $i \oplus m$ の値を加える}{in each round lane $i$ adds the value of lane $i \oplus m$ to its own}};
\end{tikzpicture}

% Width: body (120 mm); the drawing is about 104 mm wide.
\begin{tikzpicture}[x=1mm, y=1mm, font=\scriptsize\sffamily,
                    >={Stealth[length=1.2mm]}]
  \def\cw{11}
  \def\rh{11}
  \foreach \i in {0,...,7} {
    \node[text=ln-muted] at (\i*\cw+\cw/2,3*\rh+5) {\iflangja{レーン}{lane} \i};
  }
  % rows: round 0 (input), after xor 4, after xor 2, after xor 1
  \foreach \r/\lbl/\vals in {
      3/{\iflangja{入力}{input}}/{3,1,4,1,5,9,2,6},
      2/{$\oplus 4$}/{8,10,6,7,8,10,6,7},
      1/{$\oplus 2$}/{14,17,14,17,14,17,14,17},
      0/{$\oplus 1$}/{31,31,31,31,31,31,31,31}} {
    \node[anchor=east, text=ln-muted] at (-1,\r*\rh+2.5) {\lbl};
    \foreach \v [count=\i from 0] in \vals {
      \draw[fill=ln-box] (\i*\cw+1.5,\r*\rh) rectangle (\i*\cw+\cw-1.5,\r*\rh+5);
      \node at (\i*\cw+\cw/2,\r*\rh+2.5) {\v};
    }
  }
  % arrows: xor 4 (row 3 -> 2)
  \foreach \i in {0,1,2,3} {
    \draw[->, ln-accent] (\i*\cw+\cw/2,3*\rh-0.3) -- (\i*\cw+4*\cw+\cw/2,2*\rh+5.3);
    \draw[->, ln-accent] (\i*\cw+4*\cw+\cw/2,3*\rh-0.3) -- (\i*\cw+\cw/2,2*\rh+5.3);
  }
  % xor 2 (row 2 -> 1)
  \foreach \i in {0,1,4,5} {
    \draw[->, ln-accent] (\i*\cw+\cw/2,2*\rh-0.3) -- (\i*\cw+2*\cw+\cw/2,\rh+5.3);
    \draw[->, ln-accent] (\i*\cw+2*\cw+\cw/2,2*\rh-0.3) -- (\i*\cw+\cw/2,\rh+5.3);
  }
  % xor 1 (row 1 -> 0)
  \foreach \i in {0,2,4,6} {
    \draw[->, ln-accent] (\i*\cw+\cw/2,\rh-0.3) -- (\i*\cw+\cw+\cw/2,5.3);
    \draw[->, ln-accent] (\i*\cw+\cw+\cw/2,\rh-0.3) -- (\i*\cw+\cw/2,5.3);
  }
  \node[anchor=west, text=ln-muted] at (0,-4)
    {\iflangja{各段で，レーン $i$ は自分の値にレーン $i \oplus m$ の値を加える}{in each round lane $i$ adds the value of lane $i \oplus m$ to its own}};
\end{tikzpicture}

  \caption{The butterfly reduction of \cref{ex:l05-butterfly}: three
    rounds of \texttt{\_\_shfl\_xor\_sync} with masks 4, 2, 1 leave the
    sum 31 in every lane.}
  \label{fig:l05-butterfly}
\end{figure}

\begin{definition}[Warp vote function]
  A \term{warp vote function}[ワープ投票関数] evaluates a predicate $p$
  in every participating lane and combines the results across the warp
  (CPG §7.19 Warp Vote Functions):
  \begin{itemize}
    \item \texttt{\_\_all\_sync(mask, p)} is true if $p$ holds in every
      lane;
    \item \texttt{\_\_any\_sync(mask, p)} is true if it holds in at
      least one;
    \item \texttt{\_\_ballot\_sync(mask, p)} returns the 32-bit mask of
      the lanes in which it holds.
  \end{itemize}
\end{definition}

\begin{lstlisting}[language=C++, numbers=none]
unsigned m = __ballot_sync(0xffffffff, x > 0);
int count = __popc(m);              // lanes with x > 0
if (__all_sync(0xffffffff, x > 0))  // warp-uniform branch
  ...
\end{lstlisting}
A ballot followed by \texttt{\_\_popc} counts in two instructions, with
no barrier, what a shared-memory histogram would take a barrier for; a condition on the
result of \texttt{\_\_all\_sync} or \texttt{\_\_any\_sync} is uniform
across the warp by construction and never diverges.

\subsection{Cooperative groups}

\begin{definition}[Cooperative groups]
  The CUDA API that names groups of threads other than the fixed warp
  and block is called \term{cooperative groups}[Cooperative Groups] (CPG
  §8 Cooperative Groups).  A group is the whole block
  (\texttt{this\_thread\_block}), a warp or a tile of a warp
  (\texttt{tiled\_partition}), the set of lanes of a warp that are
  active at the call (\texttt{coalesced\_threads}), the whole grid, or a
  cluster.  Every
  group offers \texttt{sync()} and its size and rank; the warp-sized
  groups (tiles and coalesced groups) also offer \texttt{shfl},
  \texttt{ballot}, \texttt{any}, \texttt{all} and \texttt{cg::reduce};
  the grid group needs a cooperative launch.
\end{definition}

\Needspace{7\baselineskip}
\begin{lstlisting}[language=C++, numbers=none]
#include <cooperative_groups.h>
namespace cg = cooperative_groups;
auto block = cg::this_thread_block();
auto tile  = cg::tiled_partition<8>(block);  // 8-lane tiles
v += tile.shfl_xor(v, 4);     // shuffle within the tile
tile.sync();                  // barrier for the 8 lanes
\end{lstlisting}
The API makes the warp size and the block size parameters of a function
rather than assumptions baked into it: a routine that takes a tile as
a parameter works unchanged for tiles of 4, 8 or 32 lanes, and one that
uses only \texttt{sync()}, size and rank works for a tile and for a
block of 1024 threads alike.

\subsection{Independent thread scheduling}

\begin{definition}[Independent thread scheduling]
  The execution model introduced with Volta, called
  \term{independent thread scheduling}[独立スレッドスケジューリング],
  gives every thread of a warp its own program
  counter and call stack, so that the threads of one warp can diverge,
  interleave and reconverge at sub-warp granularity, and a thread
  waiting in a spin loop no longer blocks the other threads of its warp
  from making progress.
\end{definition}

Before Volta a warp had one program counter: on a divergent branch the
hardware ran one path to its reconvergence point before starting the
other (Lesson~6), and a thread that waited for a thread on the other
path waited for ever.  With per-thread state the scheduler may interleave
the two paths, which makes such fine-grained synchronisation within a
warp correct, at the price of giving up the guarantee that the lanes of
a warp execute in lockstep.\caution{Lecture vs document: the lesson
describes this feature as smaller warp-level operations in place of a
fixed warp size.  CPG §4.1 SIMT Architecture is more specific:
\enquote{Starting with the NVIDIA Volta architecture, Independent Thread
Scheduling allows full concurrency between threads, regardless of warp
\ldots\ the GPU maintains execution state per thread, including a
program counter and call stack}.  The warp is still 32 threads; what
changed is that code may no longer assume the lanes are in lockstep,
which is why the warp-level functions take a mask.}

\subsection{Clusters, unified on-chip memory and split barriers}

Three further additions concern the on-chip memories and
synchronisation.
\begin{description}
  \item[Shared memory across blocks.]  On Hopper the thread block
    clusters of Lesson~3 give a block access to the shared memory of the
    other blocks of its cluster, which the hardware schedules together on
    one GPU processing cluster; NVIDIA calls this \emph{distributed
    shared memory} \cite{nvidia2022hopper}.  A block can thus exchange data with
    its neighbours without a round trip through global memory.
  \item[Unified L1, shared and texture storage.]  Since Volta the L1
    data cache, the shared memory and the texture cache are one SRAM
    array per SM, split between cache and scratchpad per kernel; a
    kernel that uses little shared memory gets a large L1.  Lesson~6
    treats the carve-out and the research that led to it.\margin{H100:
    a cluster holds up to 16 blocks (a non-portable size), and an SM has
    \SI{256}{\kibi\byte} of combined L1 and shared memory, of which up
    to \SI{228}{\kibi\byte} can be shared memory
    \cite[Table~4]{nvidia2022hopper}.}
  \item[Hardware split arrive/wait barriers.]
    \begin{definition}[Split arrive/wait barrier]
      A \term{split arrive/wait barrier}[分割型到着／待機バリア] separates
      the two halves of a barrier: a thread \emph{arrives} when its data
      is ready and continues with independent work, and \emph{waits}
      only when it needs the others' data.  \texttt{\_\_syncthreads()}
      is an arrive immediately followed by a wait; the asynchronous
      barrier of CPG §7.26 (\texttt{cuda::barrier}) exposes the two
      halves, and compute capability 8.0 accelerates them in hardware.
    \end{definition}
    The gap between arrive and wait is time the thread would otherwise
    spend stalled; with the asynchronous copies of CPG §7.27, a thread
    can start the load of the next tile into shared memory, arrive, and
    compute on the current tile until the copy has landed.
\end{description}

\section{Libraries and other platforms}
\label{sec:l05-libraries}

The first rule, before any of the checklists, is to use the
libraries.\source{M5, libraries and other platforms.}

\begin{definition}[Math library]
  A \term{math library}[数学ライブラリ] is a vendor- or community-provided
  set of GPU routines for a domain (dense linear algebra, sparse
  matrices, tensors, deep-learning layers, communication) that is already
  tuned for each GPU generation: coalesced, tiled, tensor-core-aware,
  pipelined.  A program that calls one inherits those optimisations and
  follows the hardware across generations by relinking.
\end{definition}

\begin{table}[tb]
  \centering
  \small
  \caption{The NVIDIA libraries named in this lesson.}
  \label{tab:l05-libraries}
  \begin{tabular}{>{\raggedright\arraybackslash}p{20mm}
                  >{\raggedright\arraybackslash}p{86mm}}
    \toprule
    Library & Domain \\
    \midrule
    cuBLAS & dense linear algebra (BLAS: vector, matrix--vector and
      matrix--matrix operations); uses the tensor cores for GEMM \\
    cuSPARSE & sparse matrices and sparse--dense products \\
    cuTENSOR & tensor contractions, reductions and elementwise
      operations \\
    Thrust & C++ parallel algorithms on device vectors (sort, scan,
      reduce, transform), in the style of the C++ standard library \\
    NCCL & collective communication between GPUs (all-reduce,
      broadcast, all-gather), the MPI collectives of Lesson~2 for GPUs \\
    NVSHMEM & one-sided, partitioned global address space communication
      between GPUs \\
    cuDNN & deep-learning primitives: convolutions, pooling,
      normalisation, attention \\
    \bottomrule
  \end{tabular}
\end{table}

\Cref{tab:l05-libraries} lists the ones named.  The three most widely
used are \term{cuBLAS}[cuBLAS], the dense linear algebra
library whose GEMM is the kernel behind most of machine learning
(Lesson~12); \term{cuDNN}[cuDNN], the deep-learning primitives library
that frameworks such as PyTorch call; and \term{NCCL}[NCCL], the
collective-communication library for multi-GPU training (Lesson~9).
The library \term{Thrust}[Thrust] is the odd one out: not a domain library but a
C++ template library of parallel algorithms, so that a sort or a
reduction on the device is one call rather than a
kernel.\margin{\texttt{thrust::sort} on a \texttt{device\_vector} is
the quickest correct GPU sort a program can have; writing one's own only
pays when the data has structure the library cannot know
(\cref{sec:l05-sorting}).}

CUDA is NVIDIA's platform; the same ideas exist elsewhere
(\cref{tab:l05-platforms}).

\begin{table}[tb]
  \centering
  \small
  \caption{Other GPU programming platforms named in this lesson.}
  \label{tab:l05-platforms}
  \begin{tabular}{>{\raggedright\arraybackslash}p{22mm}
                  >{\raggedright\arraybackslash}p{84mm}}
    \toprule
    Platform & What it is \\
    \midrule
    OpenCL & the Khronos open standard for heterogeneous computing:
      kernels in a C dialect, portable across GPUs, CPUs and FPGAs \\
    oneAPI / SYCL & Khronos's single-source C++ model (SYCL) and Intel's
      implementation and libraries (oneAPI) \\
    OpenACC & directive-based offloading, in the style of OpenMP
      (\texttt{\#pragma acc}) \\
    HIP & AMD's C++ runtime API, source-compatible with CUDA
      (\texttt{hipMalloc}, \texttt{hipLaunchKernelGGL}); CUDA code is
      ported mechanically with \texttt{hipify} \\
    Kokkos & a C++ performance-portability layer that compiles one
      source to CUDA, HIP, SYCL or OpenMP back ends \\
    Python, Julia & high-level languages with GPU array libraries and
      kernel compilers (CuPy, Numba; CUDA.jl) \\
    \bottomrule
  \end{tabular}
\end{table}

The standard \term{OpenCL}[OpenCL] has the same execution model as CUDA under other
names (work-item, work-group, NDRange for thread, block, grid; K12
§1.3.10), \term{SYCL}[SYCL] puts host and device code in one C++ source
as CUDA does, and \term{HIP}[HIP] is close enough to CUDA that the
programs of this book translate function by function.  What transfers
between them is everything in this lesson: the three terms of
\cref{eq:l05-time} and their checklists are properties of the hardware,
not of the language.

\section{Sorting networks}
\label{sec:l05-sorting}

The sorting networks of the assigned reading \cite{kipfer2005gpusort}
are an example of an algorithm reshaped for the GPU.\source{M5, required
readings; \cite{kipfer2005gpusort}; \cite{batcher1968sorting}.}

\begin{definition}[Merge sort]
  The algorithm called \term{merge sort}[マージソート] sorts by
  recursively sorting the two halves of the sequence and \emph{merging} them: two sorted runs are
  combined by repeatedly taking the smaller of their heads.  It is
  $O(n \log n)$ and the standard sequential algorithm; but the merge is a
  data-dependent loop, which head is taken next depends on the values,
  so its control flow and its memory accesses differ from thread to
  thread.
\end{definition}

On a GPU, data-dependent control flow is divergence and data-dependent
addresses are uncoalesced accesses (\cref{tab:l05-metrics}).  The
algorithms that map well are the ones whose sequence of comparisons is
fixed in advance.

\begin{definition}[Sorting network]
  A \term{sorting network}[ソーティングネットワーク] is a fixed sequence of
  \emph{comparators}, each of which compares two positions of the array
  and swaps them if they are out of order, that sorts every input of a
  given size.  Which positions are compared at each step depends only on
  the step and the size, never on the data, so every comparator of a
  stage can execute in parallel, every thread performs the same
  operation, and the addresses are known before the kernel
  runs.\margin{Batcher introduced the two classic networks, odd-even
  merge sort and bitonic merge sort, in 1968 for hardware sorters
  \cite{batcher1968sorting}; \cite{kipfer2005gpusort} brought them to
  the GPU in 2005, before CUDA, as fragment (pixel) shader programs over
  textures (Lesson~1).}
\end{definition}

\begin{figure}[tb]
  \centering
  % A bitonic sorting network for eight inputs: six comparator stages in
% three phases (k = 2, 4, 8).  A comparator is a vertical line; the arrow
% head marks the wire that receives the larger value.  Original, after
% kipfer2005gpusort and batcher1968sorting.
% Usage: % A bitonic sorting network for eight inputs: six comparator stages in
% three phases (k = 2, 4, 8).  A comparator is a vertical line; the arrow
% head marks the wire that receives the larger value.  Original, after
% kipfer2005gpusort and batcher1968sorting.
% Usage: % A bitonic sorting network for eight inputs: six comparator stages in
% three phases (k = 2, 4, 8).  A comparator is a vertical line; the arrow
% head marks the wire that receives the larger value.  Original, after
% kipfer2005gpusort and batcher1968sorting.
% Usage: \input{figures/l05-bitonic}
% Width: body (120 mm); the drawing is about 110 mm wide.
\begin{tikzpicture}[x=1mm, y=1mm, font=\scriptsize\sffamily,
                    >={Stealth[length=1.3mm]},
                    cmp/.style={thick, ln-accent}]
  \def\wh{5.2}  % wire spacing
  \def\sx{14}   % stage spacing
  % wires
  \foreach \w in {0,...,7} {
    \draw[ln-muted] (0,-\w*\wh) -- (6*\sx+6,-\w*\wh);
    \node[anchor=east, text=ln-muted] at (-1,-\w*\wh) {$x_\w$};
  }
  % phases labels
  \node[text=ln-muted] at (\sx/2+3,5) {$k=2$};
  \node[text=ln-muted] at (2*\sx+3,5) {$k=4$};
  \node[text=ln-muted] at (4.5*\sx+3,5) {$k=8$};
  \draw[ln-muted, dashed] (\sx+3,2) -- (\sx+3,-7.5*\wh);
  \draw[ln-muted, dashed] (3*\sx+3,2) -- (3*\sx+3,-7.5*\wh);
  % a comparator from wire a to wire b (b > a); dir=1 ascending (max to b),
  % dir=0 descending (max to a)
  \newcommand{\cmpa}[3]{% x, a, b : ascending
    \draw[cmp, ->] (#1,-#2*\wh) -- (#1,-#3*\wh);
    \fill[ln-accent] (#1,-#2*\wh) circle (0.6);}
  \newcommand{\cmpd}[3]{% x, a, b : descending
    \draw[cmp, ->] (#1,-#3*\wh) -- (#1,-#2*\wh);
    \fill[ln-accent] (#1,-#3*\wh) circle (0.6);}
  % stage 1 (k=2, j=1)
  \cmpa{6}{0}{1} \cmpd{6}{2}{3} \cmpa{6}{4}{5} \cmpd{6}{6}{7}
  % stage 2 (k=4, j=2)
  \cmpa{\sx+6}{0}{2} \cmpa{\sx+9}{1}{3} \cmpd{\sx+6}{4}{6} \cmpd{\sx+9}{5}{7}
  % stage 3 (k=4, j=1)
  \cmpa{2*\sx+6}{0}{1} \cmpa{2*\sx+6}{2}{3} \cmpd{2*\sx+6}{4}{5} \cmpd{2*\sx+6}{6}{7}
  % stage 4 (k=8, j=4)
  \cmpa{3*\sx+5}{0}{4} \cmpa{3*\sx+7}{1}{5} \cmpa{3*\sx+9}{2}{6} \cmpa{3*\sx+11}{3}{7}
  % stage 5 (k=8, j=2)
  \cmpa{4*\sx+6}{0}{2} \cmpa{4*\sx+9}{1}{3} \cmpa{4*\sx+6}{4}{6} \cmpa{4*\sx+9}{5}{7}
  % stage 6 (k=8, j=1)
  \cmpa{5*\sx+6}{0}{1} \cmpa{5*\sx+6}{2}{3} \cmpa{5*\sx+6}{4}{5} \cmpa{5*\sx+6}{6}{7}
  \node[anchor=north west, text=ln-muted, align=left, text width=108mm] at (0,-7.8*\wh)
    {\iflangja{矢印の先のワイヤが大きい方の値を受け取る．$k=2$，$4$ の段は長さ $k$ の昇順・降順の列を交互に作り，$k=8$ の段が全体を昇順に並べる．}{the arrow head receives the larger value; the $k=2$ and $k=4$ phases build alternating ascending and descending runs, the $k=8$ phase sorts the whole sequence ascending}};
\end{tikzpicture}

% Width: body (120 mm); the drawing is about 110 mm wide.
\begin{tikzpicture}[x=1mm, y=1mm, font=\scriptsize\sffamily,
                    >={Stealth[length=1.3mm]},
                    cmp/.style={thick, ln-accent}]
  \def\wh{5.2}  % wire spacing
  \def\sx{14}   % stage spacing
  % wires
  \foreach \w in {0,...,7} {
    \draw[ln-muted] (0,-\w*\wh) -- (6*\sx+6,-\w*\wh);
    \node[anchor=east, text=ln-muted] at (-1,-\w*\wh) {$x_\w$};
  }
  % phases labels
  \node[text=ln-muted] at (\sx/2+3,5) {$k=2$};
  \node[text=ln-muted] at (2*\sx+3,5) {$k=4$};
  \node[text=ln-muted] at (4.5*\sx+3,5) {$k=8$};
  \draw[ln-muted, dashed] (\sx+3,2) -- (\sx+3,-7.5*\wh);
  \draw[ln-muted, dashed] (3*\sx+3,2) -- (3*\sx+3,-7.5*\wh);
  % a comparator from wire a to wire b (b > a); dir=1 ascending (max to b),
  % dir=0 descending (max to a)
  \newcommand{\cmpa}[3]{% x, a, b : ascending
    \draw[cmp, ->] (#1,-#2*\wh) -- (#1,-#3*\wh);
    \fill[ln-accent] (#1,-#2*\wh) circle (0.6);}
  \newcommand{\cmpd}[3]{% x, a, b : descending
    \draw[cmp, ->] (#1,-#3*\wh) -- (#1,-#2*\wh);
    \fill[ln-accent] (#1,-#3*\wh) circle (0.6);}
  % stage 1 (k=2, j=1)
  \cmpa{6}{0}{1} \cmpd{6}{2}{3} \cmpa{6}{4}{5} \cmpd{6}{6}{7}
  % stage 2 (k=4, j=2)
  \cmpa{\sx+6}{0}{2} \cmpa{\sx+9}{1}{3} \cmpd{\sx+6}{4}{6} \cmpd{\sx+9}{5}{7}
  % stage 3 (k=4, j=1)
  \cmpa{2*\sx+6}{0}{1} \cmpa{2*\sx+6}{2}{3} \cmpd{2*\sx+6}{4}{5} \cmpd{2*\sx+6}{6}{7}
  % stage 4 (k=8, j=4)
  \cmpa{3*\sx+5}{0}{4} \cmpa{3*\sx+7}{1}{5} \cmpa{3*\sx+9}{2}{6} \cmpa{3*\sx+11}{3}{7}
  % stage 5 (k=8, j=2)
  \cmpa{4*\sx+6}{0}{2} \cmpa{4*\sx+9}{1}{3} \cmpa{4*\sx+6}{4}{6} \cmpa{4*\sx+9}{5}{7}
  % stage 6 (k=8, j=1)
  \cmpa{5*\sx+6}{0}{1} \cmpa{5*\sx+6}{2}{3} \cmpa{5*\sx+6}{4}{5} \cmpa{5*\sx+6}{6}{7}
  \node[anchor=north west, text=ln-muted, align=left, text width=108mm] at (0,-7.8*\wh)
    {\iflangja{矢印の先のワイヤが大きい方の値を受け取る．$k=2$，$4$ の段は長さ $k$ の昇順・降順の列を交互に作り，$k=8$ の段が全体を昇順に並べる．}{the arrow head receives the larger value; the $k=2$ and $k=4$ phases build alternating ascending and descending runs, the $k=8$ phase sorts the whole sequence ascending}};
\end{tikzpicture}

% Width: body (120 mm); the drawing is about 110 mm wide.
\begin{tikzpicture}[x=1mm, y=1mm, font=\scriptsize\sffamily,
                    >={Stealth[length=1.3mm]},
                    cmp/.style={thick, ln-accent}]
  \def\wh{5.2}  % wire spacing
  \def\sx{14}   % stage spacing
  % wires
  \foreach \w in {0,...,7} {
    \draw[ln-muted] (0,-\w*\wh) -- (6*\sx+6,-\w*\wh);
    \node[anchor=east, text=ln-muted] at (-1,-\w*\wh) {$x_\w$};
  }
  % phases labels
  \node[text=ln-muted] at (\sx/2+3,5) {$k=2$};
  \node[text=ln-muted] at (2*\sx+3,5) {$k=4$};
  \node[text=ln-muted] at (4.5*\sx+3,5) {$k=8$};
  \draw[ln-muted, dashed] (\sx+3,2) -- (\sx+3,-7.5*\wh);
  \draw[ln-muted, dashed] (3*\sx+3,2) -- (3*\sx+3,-7.5*\wh);
  % a comparator from wire a to wire b (b > a); dir=1 ascending (max to b),
  % dir=0 descending (max to a)
  \newcommand{\cmpa}[3]{% x, a, b : ascending
    \draw[cmp, ->] (#1,-#2*\wh) -- (#1,-#3*\wh);
    \fill[ln-accent] (#1,-#2*\wh) circle (0.6);}
  \newcommand{\cmpd}[3]{% x, a, b : descending
    \draw[cmp, ->] (#1,-#3*\wh) -- (#1,-#2*\wh);
    \fill[ln-accent] (#1,-#3*\wh) circle (0.6);}
  % stage 1 (k=2, j=1)
  \cmpa{6}{0}{1} \cmpd{6}{2}{3} \cmpa{6}{4}{5} \cmpd{6}{6}{7}
  % stage 2 (k=4, j=2)
  \cmpa{\sx+6}{0}{2} \cmpa{\sx+9}{1}{3} \cmpd{\sx+6}{4}{6} \cmpd{\sx+9}{5}{7}
  % stage 3 (k=4, j=1)
  \cmpa{2*\sx+6}{0}{1} \cmpa{2*\sx+6}{2}{3} \cmpd{2*\sx+6}{4}{5} \cmpd{2*\sx+6}{6}{7}
  % stage 4 (k=8, j=4)
  \cmpa{3*\sx+5}{0}{4} \cmpa{3*\sx+7}{1}{5} \cmpa{3*\sx+9}{2}{6} \cmpa{3*\sx+11}{3}{7}
  % stage 5 (k=8, j=2)
  \cmpa{4*\sx+6}{0}{2} \cmpa{4*\sx+9}{1}{3} \cmpa{4*\sx+6}{4}{6} \cmpa{4*\sx+9}{5}{7}
  % stage 6 (k=8, j=1)
  \cmpa{5*\sx+6}{0}{1} \cmpa{5*\sx+6}{2}{3} \cmpa{5*\sx+6}{4}{5} \cmpa{5*\sx+6}{6}{7}
  \node[anchor=north west, text=ln-muted, align=left, text width=108mm] at (0,-7.8*\wh)
    {\iflangja{矢印の先のワイヤが大きい方の値を受け取る．$k=2$，$4$ の段は長さ $k$ の昇順・降順の列を交互に作り，$k=8$ の段が全体を昇順に並べる．}{the arrow head receives the larger value; the $k=2$ and $k=4$ phases build alternating ascending and descending runs, the $k=8$ phase sorts the whole sequence ascending}};
\end{tikzpicture}

  \caption{A bitonic sorting network for $n = 8$: six stages, four
    comparators each.  Phase $k$ merges pairs of runs of length $k/2$ into sorted
    runs of length $k$; within a phase the comparator distance is $k/2, k/4, \ldots, 1$.}
  \label{fig:l05-bitonic}
\end{figure}

\begin{definition}[Bitonic sort]
  A \emph{bitonic sequence} first rises then falls (or is a rotation of
  such a sequence).  The sorting network called
  \term{bitonic sort}[バイトニックソート] builds bitonic sequences of length $2, 4, 8, \ldots$ by
  merging sorted runs of alternating direction, and then sorts each
  bitonic sequence by a \emph{half-cleaner}: comparing element $i$ with
  element $i + k/2$ splits the sequence into two bitonic halves with
  every element of the first below every element of the second, and the
  halves are cleaned recursively (\cref{fig:l05-bitonic}).
\end{definition}

\Needspace{24\baselineskip}
For $n = 2^p$ inputs the network has $p(p+1)/2$ stages of $n/2$
comparators, i.e.\ $\frac{n}{4}\log_2 n\,(\log_2 n + 1) = O(n \log^2 n)$
comparisons, more than merge sort's $O(n \log n)$; the GPU accepts the
extra work for the regularity.  Each stage is one data-independent
pass in which thread $i$ performs the comparator at distance $j$ on its
pair.\margin{For $n = 2^{20}$:
$p = 20$, 210 stages of \num{524288} comparators, about
$1.1 \times 10^8$ comparisons, against $n \log_2 n \approx 2.1 \times
10^7$ for merge sort: five times the work, all of it regular.}  The
same idea underlies the other data-independent primitives of GPU
programming: the parallel prefix sum, the butterfly reduction of
\cref{fig:l05-butterfly}, and the FFT.

\begin{keypoint}[Lesson summary]
  A program suits a GPU when it is massively data-parallel, moves little
  data between host and device relative to its computation, and accesses
  global memory in a coalesced pattern; the profiler (Nsight Systems for
  the timeline, Nsight Compute for the kernel) finds the hotspots, on
  the CPU before porting and on the GPU after.  Execution time is
  transfer plus memory plus compute.  Transfers are shortened by pinned
  memory and hidden by asynchronous copies and a stream pipeline of
  chunks (concurrent copy and execute); zero copy and UVA let a kernel
  address host memory, but the data still crosses the bus.  Memory
  accesses are made coalesced, aligned and bank-conflict-free, and the
  effective bandwidth is compared with the theoretical one.
  Computation is reduced with cheaper instructions, lower precision,
  intrinsics, warp-uniform conditions or predication, fewer atomics and
  the tensor cores.  Warp shuffles and votes, cooperative groups,
  independent thread scheduling, distributed shared memory and split
  barriers remove old detours through shared memory and old
  restrictions of the lockstep warp.  Libraries (cuBLAS, cuDNN, NCCL,
  Thrust) carry these optimisations already, and the same rules hold on
  OpenCL, SYCL and HIP.  Sorting networks such as bitonic sort trade
  extra comparisons for a data-independent schedule, the shape a GPU
  algorithm wants.
\end{keypoint}

\end{document}
